\documentclass[oneside,openany]{final_report}
\usepackage{graphicx}
\usepackage[hidelinks]{hyperref}
\usepackage{booktabs}
\usepackage{longtable}
\usepackage{array}
\usepackage{multirow}
\usepackage{amsmath,amssymb}
\usepackage{float}
\usepackage{caption}
\usepackage{subcaption}
\usepackage{url}
\usepackage{enumitem}
\usepackage{xcolor}
\usepackage{listings}
\usepackage{tabularx}

%%%%%%%%%%%%%%%%%%%%%%
%%% Input project details
\def\studentname{Devansh Dev}
\def\reportyear{2026}
\def\projecttitle{A Novel Dual-Branch CNN with FrequencyPriorSelfAttention for Parkinson's Disease Speech Detection}
\def\headertitle{Novel Dual-Branch CNN with FP-SA for PD Speech Detection}
\def\supervisorname{Dr.\ Li Zhang}
\def\degree{BSc (Hons) in Computer Science (Artificial Intelligence)}
\def\fullOrHalfUnit{Full Unit}
\def\finalOrInterim{Final Report}

% ─── Custom Commands ───
\newcommand{\fpsa}{FrequencyPriorSelfAttention}
\newcommand{\dualcnn}{\texttt{dual\_cnn\_sa\_lstm}}

\begin{document}

\maketitle

%%%%%%%%%%%%%%%%%%%%%%
%%% Declaration

\chapter*{Declaration}

This report has been prepared on the basis of my own work. All submitted code is my original work. Where standard open-source research libraries, external tools, or published source materials have been used, these have been formally acknowledged and cited in the bibliography. Where AI tools (Claude Code, Antigravity IDE) were used for code execution and writing assistance, this is declared in Appendix~D. All architectural designs, experimental decisions, and research findings are my own original work conducted under the supervision of Dr.\ Li Zhang.

\vskip3em

Word Count: 11,065

\vskip3em

Student Name: \studentname

\vskip3em

Date of Submission: April 4, 2026

\vskip3em

Signature: \includegraphics[height=1.2cm]{figures/signature.png}

%%%%%%%%%%%%%%%%%%%%%%
%%% Table of Contents
\newpage
\tableofcontents\pdfbookmark[0]{Table of Contents}{toc}\newpage

%%%%%%%%%%%%%%%%%%%%%%
%%% Abstract

\begin{abstract}

Parkinson's Disease (PD) affects over ten million people worldwide, and speech degradation is among the earliest and most consistently observed non-invasive biomarkers of the condition. Existing deep learning approaches to PD speech detection frequently rely on single-dataset evaluation, lack formal data leakage controls, and do not assess cross-lingual transferability. This thesis presents a novel \dualcnn{} architecture that fuses dual-branch convolutional neural networks, specifically ResNet50 and EfficientNetV2-S, with a custom \fpsa{} (FP-SA) module, Coordinate Attention, and a Bidirectional Long Short-Term Memory (BiLSTM) temporal head for binary PD versus healthy control classification from log-Mel spectrogram representations. The system is evaluated under speaker-independent five-fold cross-validation with a certified zero-leakage guard, producing $N\!=\!15$ observations per dataset (three seeds $\times$ five folds) and employing Wilcoxon signed-rank non-parametric testing for statistical rigour.

On the Italian PD dataset, the proposed architecture achieves a mean macro F1 of $0.964 \pm 0.048$ ($N\!=\!15$) under speaker-independent cross-validation with certified zero data leakage. To the best of the author's knowledge, this is the strongest result reported on this dataset under a fully verifiable evaluation protocol. Cross-lingual evaluation on the Colombian Spanish PC-GITA diadochokinetic corpus yields F1\,=\,0.780, quantifying a meaningful generalisation gap across languages and speech tasks. A systematic fold-level analysis reveals a persistent performance anomaly on Fold~3 of the Italian PD dataset (consistent across all architectures and all random seeds), traced to approximately 21 healthy control subjects whose speech spectrograms consistently exhibit PD-like spectral patterns. This subgroup does not appear to have been previously identified in the literature.

The principal contributions are: (1)~the \dualcnn{} architecture with a novel FP-SA frequency-prior attention mechanism; (2)~a certified zero-leakage evaluation pipeline with an active RuntimeError guard; (3)~cross-lingual PD evaluation under a unified experimental framework; (4)~identification and characterisation of the Fold~3 borderline healthy control cluster; and (5)~an evidenced six-phase architectural progression from baseline CNN to the final system.

\end{abstract}
\newpage

%%%%%%%%%%%%%%%%%%%%%%
%%% Project Specification

\chapter*{Project Specification}
\addcontentsline{toc}{chapter}{Project Specification}

\textbf{Title:} A Novel Dual-Branch CNN with FrequencyPriorSelfAttention for Parkinson's Disease Speech Detection

\textbf{Supervisor:} Dr.\ Li Zhang

\textbf{Objectives:}
\begin{enumerate}
    \item Design and implement a novel dual-branch CNN architecture (\dualcnn{}) combining ResNet50 and EfficientNetV2-S with a custom FrequencyPriorSelfAttention module for Parkinson's Disease speech detection from log-Mel spectrograms.
    \item Establish a certified zero-leakage evaluation pipeline with speaker-independent five-fold cross-validation and an active RuntimeError guard.
    \item Evaluate the system across four datasets (Italian PD, PC-GITA, ESC-50, EmoDB) with $N\!=\!15$ statistical validation per dataset.
    \item Conduct a controlled attention ablation study isolating the contribution of each architectural component.
    \item Perform cross-lingual PD evaluation (Italian and Colombian Spanish) under identical experimental conditions.
    \item Compare against published state-of-the-art results with honest statistical analysis.
\end{enumerate}

\textbf{Deliverables:}
\begin{itemize}
    \item Complete codebase: model architectures, training engine, spectrogram generation pipeline, leakage verification tools, and Kaggle execution notebooks.
    \item Experimental results: 75+ cross-validated training runs with JSON summaries, confusion matrices, and TensorBoard logs.
    \item Final report documenting methodology, architecture, results, and analysis.
    \item Conference paper submitted for publication.
\end{itemize}

% ═══════════════════════════════════════════════════════════════
% CHAPTER 1 — INTRODUCTION
% ═══════════════════════════════════════════════════════════════
\chapter{Introduction}

\section{Clinical Motivation}

Parkinson's Disease is a progressive neurodegenerative disorder affecting the dopaminergic neurons of the substantia nigra, with global prevalence exceeding ten million cases and a projected doubling by 2040~\cite{dorsey2018}. Among the earliest and most pervasive symptoms are speech and voice impairments: approximately 90\% of PD patients develop some form of dysarthria during the course of the disease~\cite{ho1999}, manifesting as reduced loudness, monotone pitch, imprecise articulation, and breathy voice quality~\cite{duffy2013}. These vocal biomarkers emerge years before motor symptoms reach clinical detection thresholds, positioning speech analysis as a non-invasive, low-cost screening modality with substantial clinical potential.

Traditional PD diagnosis relies on subjective clinical assessment of motor symptoms using rating scales such as the UPDRS, which suffer from inter-rater variability and limited sensitivity to early-stage disease. Acoustic analysis offers an objective, repeatable, and scalable alternative. However, translating this potential into reliable automated detection systems requires addressing fundamental challenges in machine learning methodology that the current literature has not adequately resolved.

\section{Research Gap}

Despite significant progress in applying deep learning to PD speech detection, the field exhibits three systematic weaknesses that undermine confidence in published results.

First, \textbf{data leakage} remains poorly controlled. Many studies employ random train-test splits on datasets where multiple recordings originate from the same speaker, allowing the model to memorise speaker identity rather than learn pathological speech features. Formal leakage verification is rarely reported, and when it is, the methodology is seldom reproducible~\cite{sakar2013, tsanas2012}.

Second, \textbf{single-dataset evaluation} is the norm. The majority of published PD speech detection systems are trained and tested on a single corpus, providing no evidence of cross-lingual or cross-task generalisation. Given the well-documented influence of language-specific phonological systems on PD speech manifestation~\cite{orozco2016}, this omission represents a critical gap.

Third, \textbf{statistical reporting} is frequently inadequate. Single-split evaluations with unreported variance provide no basis for statistical comparison, yet such results are routinely positioned as state-of-the-art claims.

\section{Research Questions}

This thesis addresses four research questions that collectively target the identified gaps:

\begin{enumerate}[label=\textbf{RQ\arabic*:}]
    \item Can a dual-branch CNN architecture with frequency-aware self-attention achieve state-of-the-art PD speech detection performance under speaker-independent cross-validation with formal leakage certification?
    \item How does the proposed FrequencyPriorSelfAttention module compare against standard attention mechanisms (Coordinate Attention, Attention Gate, Squeeze-and-Excitation) when evaluated as the sole experimental variable on clinical speech data?
    \item To what extent do PD speech detection models trained on one language (Italian) generalise to a different language and speech task (Colombian Spanish diadochokinetic analysis)?
    \item What does systematic fold-level instability analysis reveal about the composition of clinical speech datasets, and what are the clinical implications of identified anomalies?
\end{enumerate}

\section{Contributions}

The principal contributions of this work are:

\begin{enumerate}
    \item \textbf{Novel Architecture.} The \dualcnn{} model, a dual-branch architecture combining ResNet50 and EfficientNetV2-S feature extractors with a custom FrequencyPriorSelfAttention module that injects a learnable frequency-axis prior bias into the self-attention mechanism, designed specifically for the spectral structure of PD dysarthria.
    \item \textbf{Zero-Leakage Certification.} A formal data leakage verification pipeline comprising \texttt{StratifiedGroupKFold} splitting by subject identifier, an independent audit script (\texttt{verify\_leakage.py}), and an inline RuntimeError guard in the training engine that halts execution upon any detected subject overlap. This evaluation protocol exceeds the reproducibility standards of most published work on the Italian PD dataset.
    \item \textbf{Cross-Lingual PD Evaluation.} The first evaluation, to the author's knowledge, of a unified PD detection pipeline across both the Italian PD corpus and the Colombian Spanish PC-GITA diadochokinetic corpus under identical experimental conditions.
    \item \textbf{Fold~3 Clinical Anomaly.} Identification and characterisation of a borderline healthy control cluster in the Italian PD dataset whose speech spectrograms consistently mimic PD patterns across all architectures and all random seeds, a finding with implications for clinical screening deployment.
    \item \textbf{Evidenced Architectural Progression.} A six-phase experimental trajectory in which every design decision is motivated by empirical evidence from the preceding phase, documented through 75+ cross-validated training runs.
    \item \textbf{Dissemination.} The core findings of this thesis have been submitted for publication as an IEEE conference paper. The paper presents the \dualcnn{} architecture, the zero-leakage evaluation protocol, and the cross-lingual generalisation results in condensed form for peer review.
\end{enumerate}

\section{Advancement Beyond Interim Report}
\label{sec:advancement}

The interim report submitted in Term~1 (scoring 93/100) presented a comparative study of CNN architectures with Squeeze-and-Excitation and CBAM attention modules on two non-clinical audio datasets (ESC-50 environmental sound classification and EmoDB speech emotion recognition). The final project represents a significant shift in research scope, methodology, and technical contribution. Table~\ref{tab:advancement} quantifies this progression across nine dimensions.

\begin{table}[H]
\centering
\caption{Advancement from Interim Report to Final Project}
\label{tab:advancement}
\small
\begin{tabularx}{\textwidth}{l X X}
\toprule
\textbf{Dimension} & \textbf{Interim Report} & \textbf{Final Report} \\
\midrule
Clinical focus & None (env.\ sounds + emotion) & PD speech detection (clinical biomarker) \\
Datasets & 2 (ESC-50, EmoDB) & 4 (Italian PD, PC-GITA, ESC-50, EmoDB) \\
Evaluation protocol & 80/20 single split & Speaker-independent 5-fold CV \\
Statistical rigour & 3 seeds & $N\!=\!15$ (3 seeds $\times$ 5 folds), Wilcoxon testing \\
Architecture & Standard CNN variants & Novel \dualcnn{} with FP-SA \\
Leakage control & Grouped split fix & Formal certification + inline RuntimeError guard \\
Cross-lingual evaluation & None & Italian PD + Colombian Spanish PC-GITA \\
Clinical finding & None & Fold~3 anomaly: hard HC borderline cluster \\
SOTA comparison & None & Strongest verifiable result on Italian PD (F1\,=\,0.964); no prior work matches the evaluation rigour \\
\bottomrule
\end{tabularx}
\end{table}

Across all nine dimensions, the final project goes well beyond what was submitted at interim. The research question was reframed, a novel architecture was introduced, and a stricter evaluation methodology was applied than is typical in the published literature on these datasets.

\section{Report Structure}

The remainder of this thesis is organised as follows. Chapter~2 presents the theoretical background underpinning the system. Chapter~3 reviews and critically evaluates the relevant literature. Chapter~4 describes the experimental methodology, datasets, and reproducibility guarantees. Chapter~5 details the system architecture and its evolution across six experimental phases. Chapter~6 covers implementation specifics. Chapter~7 presents results with theoretical analysis. Chapter~8 provides critical discussion including threats to validity and failure case analysis. Chapter~9 concludes the thesis. Chapter~10 addresses professional and ethical issues. Chapter~11 presents the project schedule, and Chapter~12 summarises the project diary.

% ═══════════════════════════════════════════════════════════════
% CHAPTER 2 — BACKGROUND THEORY
% ═══════════════════════════════════════════════════════════════
\chapter{Background Theory}

This chapter establishes the theoretical foundations required to understand the architectural decisions and experimental methodology of this thesis. Each concept is presented with sufficient depth to justify its role in the system design, with particular attention to elements that were not covered in the interim report.

\section{Audio Representation: From Waveform to Spectrogram}

Raw audio waveforms encode amplitude over time but lack explicit frequency information. The Short-Time Fourier Transform (STFT) addresses this by applying the Discrete Fourier Transform to overlapping windowed segments of the signal, producing a complex-valued time-frequency representation. For a discrete signal $x[n]$ with window function $w[n]$ of length $N$:

\begin{equation}
    X[m, k] = \sum_{n=0}^{N-1} x[n + mH] \cdot w[n] \cdot e^{-j2\pi kn / N}
    \label{eq:stft}
\end{equation}

\noindent where $m$ indexes the time frame, $k$ indexes the frequency bin, and $H$ is the hop length in samples.

The magnitude spectrogram $|X[m,k]|$ captures energy distribution across time and frequency but allocates linear resolution to the frequency axis, which poorly matches human auditory perception, which is approximately logarithmic above 1\,kHz. The Mel-frequency filterbank addresses this by applying a set of $M$ triangular filters spaced according to the Mel scale:

\begin{equation}
    S_{\text{mel}}[m, i] = \sum_{k} |X[m,k]|^2 \cdot H_i[k], \quad i = 1, \ldots, M
    \label{eq:mel}
\end{equation}

\noindent where $H_i[k]$ is the $i$-th triangular filter. Applying a logarithmic compression yields the log-Mel spectrogram, which compresses dynamic range and approximates the nonlinear loudness perception of the human auditory system. In this work, all spectrograms are generated with $n_{\text{fft}} = 2048$, hop length $= 512$, and $M = 128$ Mel bands at a target sample rate of 16\,kHz, producing representations that are resized to $224 \times 224$ pixels for CNN ingestion.

\section{Convolutional Neural Network Architectures}

\subsection{ResNet50: Residual Learning}

He et al.~\cite{he2016} introduced residual connections that enable training of substantially deeper networks by learning residual mappings $\mathcal{F}(\mathbf{x}) + \mathbf{x}$ rather than direct transformations. ResNet50 comprises four residual stages with bottleneck blocks (1$\times$1, 3$\times$3, 1$\times$1 convolutions), producing a 2048-dimensional feature map at the final stage. Its identity shortcuts prevent gradient degradation, making it a robust backbone for transfer learning from ImageNet to spectrogram domains.

In this work, ResNet50 serves as the primary feature extractor in both the single-branch baseline and as one branch of the dual-branch architecture, with ImageNet-pretrained weights (V2) providing initialisation.

\subsection{EfficientNetV2-S: Compound Scaling}

Tan and Le~\cite{tan2021} proposed EfficientNetV2, which jointly optimises network depth, width, and resolution through compound scaling coefficients. The ``S'' (small) variant employs Fused-MBConv blocks in early stages and standard MBConv blocks with squeeze-and-excitation in later stages, achieving superior accuracy-to-FLOPS efficiency compared to ResNet variants.

EfficientNetV2-S produces 1280-dimensional feature maps at its final convolutional layer. Its compound scaling captures fine-grained frequency details that complement ResNet50's structural features, providing the theoretical basis for the dual-branch design. This architectural choice was preferred over VGG16 (excessive parameter count, no residual connections), MobileNetV2 (empirically weaker on clinical datasets in Phase~2 evaluation), and DenseNet (diminishing returns from dense connectivity on small clinical datasets).

\section{Bidirectional LSTM for Temporal Modelling}

Long Short-Term Memory (LSTM) networks~\cite{hochreiter1997} address the vanishing gradient problem in recurrent architectures through gated memory cells that selectively retain, update, or discard information across time steps. The Bidirectional extension~\cite{graves2005} processes the input sequence in both forward and reverse directions, capturing both onset-to-peak and decay patterns.

For PD speech spectrograms, pathological tremor manifests as temporally persistent disruptions in formant structure. Standard CNN architectures with global average pooling collapse the temporal axis entirely, suppressing frame-to-frame pathological patterns through translational invariance. A BiLSTM operating on the preserved time columns of the feature map explicitly models these temporal dependencies. In the \dualcnn{} architecture, the spatial grid ($7 \times 7$) is reshaped into a 49-step sequence, allowing the BiLSTM to capture cross-frame relationships that spatial pooling would destroy.

\section{FrequencyPriorSelfAttention (FP-SA)}
\label{sec:fpsa_theory}

Standard multi-head self-attention~\cite{vaswani2017} computes pairwise token affinities via scaled dot-product attention:

\begin{equation}
    \text{Attention}(\mathbf{Q}, \mathbf{K}, \mathbf{V}) = \text{softmax}\!\left(\frac{\mathbf{Q}\mathbf{K}^\top}{\sqrt{d_k}}\right) \mathbf{V}
    \label{eq:standard_attn}
\end{equation}

This formulation treats all token pairs identically regardless of their spectral distance. In PD dysarthria, pathological features manifest as structured disruptions at specific harmonic frequency bands, such as vocal tremor at 4--8\,Hz modulation and formant F0 shifts in the 70--300\,Hz range~\cite{rusz2011}. An isotropic attention mechanism ignores this physical structure.

FrequencyPriorSelfAttention introduces a frequency-axis prior bias $\mathbf{B}_{\text{freq}}$ into the attention computation:

\begin{equation}
    \text{FP-SA}(\mathbf{Q}, \mathbf{K}, \mathbf{V}) = \text{softmax}\!\left(\frac{\mathbf{Q}\mathbf{K}^\top + \mathbf{B}_{\text{freq}}}{\sqrt{d_k}}\right) \mathbf{V}
    \label{eq:fpsa}
\end{equation}

\noindent where:

\begin{itemize}[nosep]
    \item $\mathbf{Q} \in \mathbb{R}^{T \times d_k}$: query matrix, $T = H \times W$ is the number of time-frequency tokens, $d_k$ is the key dimension
    \item $\mathbf{K} \in \mathbb{R}^{T \times d_k}$: key matrix
    \item $\mathbf{V} \in \mathbb{R}^{T \times d_v}$: value matrix, $d_v$ is the value dimension
    \item $\mathbf{B}_{\text{freq}} \in \mathbb{R}^{T \times T}$: frequency-axis prior bias matrix
\end{itemize}

The bias encodes the inductive prior that tokens from nearby frequency bins are more likely to be semantically related:

\begin{equation}
    \mathbf{B}_{\text{freq}}[i,j] = \beta \cdot |\text{freq\_bin}(i) - \text{freq\_bin}(j)|^{-1}
    \label{eq:freq_bias}
\end{equation}

\noindent where $\text{freq\_bin}(i)$ returns the frequency bin index of token $i$, and $\beta$ is a learnable scalar controlling the prior strength.

The prior is grounded in the acoustic properties of PD speech pathology: harmonically related frequency components in a spectrogram carry correlated diagnostic information. When $\mathbf{B}_{\text{freq}} = \mathbf{0}$, FP-SA reduces to standard self-attention, making the prior an additive inductive bias rather than a hard constraint. To the best of our knowledge, this is the first application of a frequency-indexed self-attention prior to Parkinson's Disease speech detection.

In implementation, the frequency prior is injected into the key computation rather than directly into the attention score matrix: a learnable embedding $\mathbf{F} \in \mathbb{R}^{H \times d}$ (one vector per frequency row) is added to the key projections before computing $\mathbf{Q}\mathbf{K}^\top$. This is mathematically equivalent to Equation~\ref{eq:fpsa} but computationally more efficient, as it avoids materialising the full $T \times T$ bias matrix.

\section{Coordinate Attention}

Hou et al.~\cite{hou2021} proposed Coordinate Attention (CA), which decomposes spatial attention into two orthogonal one-dimensional pooling operations, horizontal (time-axis) and vertical (frequency-axis), before recombining through learned channel weights. Unlike Squeeze-and-Excitation~\cite{hu2018}, which collapses spatial dimensions entirely, CA preserves directional information, enabling the network to attend to frequency-specific or time-specific patterns independently.

For spectrogram processing, this directional decomposition is theoretically advantageous: PD-related formant shifts are primarily frequency-axis phenomena, while speech timing disruptions are time-axis phenomena. CA can learn to weight these axes independently. In the \dualcnn{} architecture, CA is applied after FP-SA to provide direction-aware feature recalibration on the attention-enhanced feature maps.

\section{Attention Gate}

Originally developed for medical image segmentation, the Attention Gate uses a global-average-pooled gating signal to learn a spatial suppression mask that filters irrelevant regions. In the context of this work, it was evaluated as an alternative to CA during Phase~4 (attention ablation). While theoretically appealing for its spatial selectivity, empirical evaluation revealed that it adds approximately 8M parameters (compared to $\sim$2K for CA at layer4) and did not improve over CA on the clinical datasets, leading to its exclusion from the final architecture.

\section{Transfer Learning: ImageNet to Clinical Spectrograms}

Transfer learning from ImageNet-pretrained models to spectrogram classification builds on the well-documented finding~\cite{he2016} that early convolutional layers learn universal visual features (edges, textures, colour gradients) that transfer effectively across visual domains. Log-Mel spectrograms, when rendered as three-channel RGB images, share sufficient structural similarity with natural images for this transfer to be effective.

However, the domain gap between ImageNet and clinical spectrograms is significant. ImageNet features are optimised for object recognition with diverse spatial configurations, whereas spectrograms have a fixed axis semantics (frequency vertical, time horizontal) and pathological signals are often subtle texture variations rather than salient objects. This domain gap motivates: (a)~fine-tuning deep layers rather than using fixed features; (b)~using static ImageNet normalisation parameters to avoid leaking training-set statistics; and (c)~designing frequency-aware attention modules that encode spectrogram-specific structure.

\section{Speaker-Independent K-Fold Cross-Validation}

Standard $K$-fold cross-validation assigns individual samples to folds. In clinical speech datasets where each subject contributes multiple recordings, this permits the model to encounter different recordings from the same speaker in both training and validation sets, a form of data leakage that inflates performance estimates~\cite{sakar2013}.

Speaker-independent $K$-fold CV enforces a group constraint: all recordings from a given subject appear in exactly one fold. \texttt{StratifiedGroupKFold} from scikit-learn~\cite{sklearn2011} implements this by treating the subject identifier as the group key while maintaining approximate class balance across folds. With $K\!=\!5$ and three random seeds, this produces $N\!=\!15$ observations per dataset, sufficient for non-parametric statistical testing.

\section{Wilcoxon Signed-Rank Test}

The Wilcoxon signed-rank test is a non-parametric paired difference test that does not assume normality of the underlying distribution. For $K$-fold CV comparisons with $K\!=\!5$ paired observations, the normality assumption required by the paired $t$-test cannot be reliably verified. The Wilcoxon test ranks the absolute differences between paired fold-level scores and tests whether the median difference is zero. This provides a statistically conservative comparison between architectural variants without distributional assumptions.

\section{SpecAugment}

Park et al.~\cite{park2019} introduced SpecAugment, a data augmentation technique that applies random time and frequency masking to spectrogram inputs during training. Given a spectrogram input, a contiguous block of $f$ frequency rows (where $f \sim \text{Uniform}(0, F)$) is zeroed, and independently, a contiguous block of $t$ time columns (where $t \sim \text{Uniform}(0, T)$) is zeroed. This regularisation is domain-appropriate for clinical speech because it simulates the natural variability of speech signals, including silence gaps (time masking) and frequency-selective noise (frequency masking), without modifying the underlying waveform. In this work, $T = F = 40$ pixels on the $224 \times 224$ input.

% ═══════════════════════════════════════════════════════════════
% CHAPTER 3 — LITERATURE REVIEW
% ═══════════════════════════════════════════════════════════════
\chapter{Literature Review}

This chapter critically evaluates existing work across five research areas relevant to this thesis: PD speech detection, attention mechanisms in audio classification, transfer learning for medical audio, cross-lingual speech pathology, and data leakage in clinical machine learning.

\section{PD Speech Detection}

Sakar et al.~\cite{sakar2013} established an early benchmark for PD speech detection using traditional machine learning features (jitter, shimmer, harmonic-to-noise ratio) extracted from sustained vowel phonations, achieving classification accuracies in the 75--85\% range. While foundational, their feature engineering approach is inherently limited by the manually selected acoustic parameters and cannot capture the full spectro-temporal complexity of PD dysarthria. Tsanas et al.~\cite{tsanas2012} extended this work with nonlinear dynamics features and achieved higher performance on a proprietary dataset, but their evaluation used subject-dependent splitting, raising concerns about leakage-inflated results.

Aversano et al.~\cite{aversano2024} report F1\,=\,0.971 on the Italian PD dataset using a machine learning pipeline with acoustic feature extraction, but do not disclose whether speaker-independent splits were used or whether any data leakage controls were applied. Without knowing the splitting strategy, it is unclear whether their reported figure reflects genuine generalisation or benefits from speaker-level leakage. This thesis addresses that gap by evaluating under a fully documented, speaker-independent protocol with certified zero leakage, producing results that are directly reproducible and verifiable.

Rusz et al.~\cite{rusz2011} demonstrated that quantitative acoustic measurements, particularly in formant transitions, voice onset time, and speech rate, can characterise early untreated PD with clinical sensitivity. Their work establishes the acoustic basis for the FP-SA frequency prior: pathological features cluster at specific frequency bands rather than being uniformly distributed, justifying a non-isotropic attention mechanism.

\section{Attention Mechanisms in Audio Classification}

Hu et al.~\cite{hu2018} introduced Squeeze-and-Excitation (SE) networks, which recalibrate channel-wise feature responses using global average pooling followed by a bottleneck excitation. While effective for ImageNet classification, SE collapses all spatial information into a single vector, discarding the frequency-time structure critical for spectrogram interpretation. Experimental evaluation in Term~1 confirmed that SE provides marginal gains on ESC-50 ($+0.1\%$ F1) but negligible improvement on EmoDB, consistent with the loss of spatial specificity.

Woo et al.~\cite{woo2018} extended SE with spatial attention in CBAM, applying sequential channel and spatial attention modules. However, Term~1 experiments revealed severe training instability with CBAM on both ESC-50 (F1 variance $\pm 10.9\%$) and EmoDB (F1 variance $\pm 19.3\%$). This instability was attributed to the sequential channel-then-spatial attention pipeline amplifying noise in the spatial attention branch when operating on spectrogram inputs with relatively uniform spatial structure.

Coordinate Attention~\cite{hou2021} addresses the limitations of both SE and CBAM by decomposing attention into direction-aware components without collapsing spatial dimensions entirely. Its factorised design, with horizontal pooling for the time-axis and vertical pooling for the frequency-axis, maps naturally onto the axes of a spectrogram, providing the directional specificity that SE lacks and the stability that CBAM fails to deliver. Phase~4 ablation confirmed CA as the strongest single-module attention mechanism for clinical speech data, motivating its inclusion in the final architecture.

\section{Transfer Learning for Medical Audio}

The efficacy of ImageNet transfer learning for audio classification is well established for non-clinical domains. However, clinical speech pathology presents a more challenging domain gap: pathological vocal features are subtle spectral texture variations rather than the salient objects that ImageNet features are optimised to detect. Dimauro et al.~\cite{dimauro2019} provided the Italian PD corpus but did not employ deep transfer learning, relying instead on classical features. The gap between classical and deep approaches on this dataset motivates the systematic comparison undertaken in this thesis.

A critical consideration is normalisation: computing training-set mean and standard deviation for input normalisation, as is standard practice, introduces a subtle form of data leakage when the same statistics are applied at inference. This work eliminates this risk by using fixed ImageNet normalisation parameters ($\mu = [0.485, 0.456, 0.406]$, $\sigma = [0.229, 0.224, 0.225]$) throughout, ensuring that no training-set statistics propagate to the validation forward pass.

\section{Cross-Lingual Speech Pathology}

Orozco-Arroyave et al.~\cite{orozco2014} created the PC-GITA corpus, a Colombian Spanish dataset including diadochokinetic (DDK) tasks, sustained vowels, and read speech from 50 PD and 51 HC subjects. Their subsequent cross-lingual analysis~\cite{orozco2016} demonstrated that automatic PD detection accuracy varies significantly across languages, attributing this to language-specific phonological constraints on the manifestation of dysarthria.

No prior work, to the author's knowledge, has evaluated a single deep learning pipeline on both the Italian PD corpus and the PC-GITA DDK corpus under identical experimental conditions (same architecture, same hyperparameters, same evaluation protocol). This thesis fills that gap, providing the first direct measurement of cross-lingual generalisation for a spectrogram-based PD detection system.

\section{Data Leakage in Clinical Machine Learning}

Data leakage, the inadvertent inclusion of test-time information in the training process, is a pervasive problem in clinical machine learning that is increasingly recognised as a threat to published results. The specific form relevant to clinical speech is \emph{subject leakage}: when multiple recordings from the same patient appear in both training and test sets, the model can achieve artificially high performance by recognising speaker identity rather than pathological features.

Two independent leakage events were discovered and corrected during this project. The first, in Term~1, involved augmentation leakage: augmented variants of training recordings were not being grouped with their source clips during splitting. The fix (grouped splitting on source clip ID) was the genesis of the formal leakage certification pipeline developed in Term~2. The second, discovered during Phase~5, was a Kaggle path-resolution bug: stale fold CSVs from a prior dataset version caused Folds~0 and 1 to load corrupted splits, producing bimodal F1 distributions (perfect 1.0 on corrupted folds, realistic 0.85--0.93 on correct folds). This was fixed by hardcoding the Kaggle dataset slug and implementing an inline RuntimeError leakage guard.

Both bugs were caught and corrected before any final results were generated. The experience shows that even with careful experimental design, leakage can be introduced through non-obvious mechanisms, which is precisely why active verification is necessary rather than relying on split descriptions alone.

% ═══════════════════════════════════════════════════════════════
% CHAPTER 4 — METHODOLOGY
% ═══════════════════════════════════════════════════════════════
\chapter{Methodology}

\section{Datasets}

Four datasets spanning clinical and non-clinical audio domains are employed. Table~\ref{tab:datasets} summarises their characteristics.

\begin{table}[H]
\centering
\caption{Dataset Summary}
\label{tab:datasets}
\small
\begin{tabularx}{\textwidth}{l l l c c X}
\toprule
\textbf{Dataset} & \textbf{Language} & \textbf{Task} & \textbf{Subjects} & \textbf{Files} & \textbf{Split Strategy} \\
\midrule
Italian PD & Italian & PD vs HC & $\sim$65 & 831 & Speaker-ind.\ 5-fold \\
PC-GITA & Col.\ Spanish & DDK analysis & 101 & 600 & Speaker-ind.\ 5-fold \\
ESC-50 & N/A & 50-class env.\ sounds & 2000 clips & $\sim$8k & Grouped 5-fold \\
EmoDB & German & 7-class emotion & 10 actors & $\sim$3.2k & Speaker-ind.\ 5-fold \\
\bottomrule
\end{tabularx}
\end{table}

The \textbf{Italian PD} corpus~\cite{dimauro2019} contains sustained vowels and read speech from approximately 65 subjects (PD and age-matched healthy controls). It serves as the primary evaluation dataset for ablation studies and SOTA comparison. The \textbf{PC-GITA} corpus~\cite{orozco2014} comprises DDK recordings (syllable repetition tasks: \textit{pa-pa-pa}, \textit{ta-ta-ta}, \textit{ka-ka-ka}, \textit{pataka}, \textit{pakata}, \textit{petaka}) from 101 speakers (51 HC, 50 PD). DDK tasks are the clinical gold standard for motor speech assessment in PD, as they stress the articulatory subsystem more directly than read speech. PC-GITA serves exclusively as the cross-lingual generalisation benchmark, not for ablation.

\textbf{ESC-50}~\cite{piczak2015} (50-class environmental sound classification, 2000 clips) and \textbf{EmoDB}~\cite{burkhardt2005} (7-class German speech emotion, 10 actors, 535 utterances) serve as non-clinical reference datasets. They provide evidence that the system is not overfit to clinical speech and establish baseline performance benchmarks against which the clinical-specific design choices can be evaluated.

\section{Spectrogram Generation}

All datasets are processed through a unified spectrogram pipeline. Raw audio is resampled to 16\,kHz (Italian PD, PC-GITA, EmoDB, ESC-50). The log-Mel spectrogram is computed with $n_{\text{fft}} = 2048$, hop length $= 512$, and $n_{\text{mels}} = 128$, then saved as $224 \times 224$ RGB PNG images. Clinical datasets (Italian PD, PC-GITA) use no audio-level augmentation to preserve subtle pathological features. ESC-50 training spectrograms include noise injection, pitch shift ($+2$ semitones), and time stretch ($0.9\times$) to increase data diversity for the 50-class task. SpecAugment (time and frequency masking with $T = F = 40$ pixels) is applied at training time for Phase~5 and Phase~6 experiments on clinical datasets.

\begin{figure}[htbp]
    \centering
    \includegraphics[width=0.95\linewidth]{figures/fig_spec_grid.png}
    \caption{Representative log-Mel spectrograms from each dataset used in this study (128 Mel bands, $n_{\text{fft}}=2048$, hop length $=512$, 16\,kHz). \textit{Top-left:} Italian PD --- sustained vowel from a PD patient, showing characteristic harmonic blurring and reduced high-frequency energy. \textit{Top-right:} PC-GITA DDK (/ka/ repetition) --- rhythmic articulatory bursts with irregular inter-burst timing typical of hypokinetic dysarthria. \textit{Bottom-left:} ESC-50 (dog bark) --- broadband noise with transient tonal structure, illustrating the diversity of the non-clinical reference domain. \textit{Bottom-right:} EmoDB --- emotional speech showing clear formant structure and dynamic energy variation. All spectrograms are rendered as $224\times224$ RGB images prior to CNN ingestion.}
    \label{fig:spec_grid}
\end{figure}

\section{Zero-Leakage Certification Protocol}

The leakage certification protocol operates at three levels.

\textbf{Level 1: Grouped Splitting.} Clinical datasets are split using \texttt{StratifiedGroupKFold} with \texttt{subject\_id} as the group key. This guarantees that all recordings from a given subject appear in exactly one fold. Non-clinical datasets use clip-level grouping to prevent augmented variants from crossing the train-validation boundary.

\textbf{Level 2: Independent Audit.} The \texttt{verify\_leakage.py} script independently loads all train and validation CSVs and verifies: (a)~zero file path overlap between train and validation sets; (b)~zero subject overlap for clinical datasets; (c)~correct class distribution across folds. This audit was executed across all 50 fold CSV files ($5 \text{ datasets} \times 5 \text{ folds} \times 2 \text{ (train/val)}$) with zero violations detected.

\textbf{Level 3: Active Runtime Guard.} An inline assertion in \texttt{train\_unified.py} reads the train and validation CSVs at the start of every training run, computes the subject-ID intersection, and raises a \texttt{RuntimeError} with an explicit data leakage message if any overlap is detected. This guard executes before any model training begins, providing an active defence that would catch leakage introduced by environment-specific path resolution (as occurred during Phase~5 on Kaggle).

The leakage certification protocol implemented here, including an active RuntimeError guard that halts execution on any detected subject overlap, provides a stronger guarantee than the split descriptions reported in most comparable published works.

\section{Validation-as-Holdout Justification}

This work employs a two-way split (train/validation) rather than a three-way split (train/validation/test). The rationale is threefold. First, clinical datasets have limited subjects ($\sim$65 for Italian PD, 101 for PC-GITA); a tertiary test split would reduce either training capacity or produce a statistically unreliable test set. Second, speaker-independent K-fold CV already provides an unbiased generalisation estimate, since each fold's validation set contains subjects never seen during training. Third, the $N\!=\!15$ multi-seed evaluation (3 seeds $\times$ 5 folds) provides substantially more statistical power than a single evaluation on a small held-out test set.

\section{Tiered Regularisation Protocol}

A universal regularisation configuration failed catastrophically during Phase~1 (the ``regularisation pendulum'': ESC-50 and EmoDB underfitted with aggressive clinical parameters, while Pitt Corpus overfitted with standard parameters). This motivated a dataset-specific tiered protocol, shown in Table~\ref{tab:tiers}.

\begin{table}[H]
\centering
\caption{Tiered Regularisation Protocol}
\label{tab:tiers}
\begin{tabular}{l l c c c c}
\toprule
\textbf{Tier} & \textbf{Datasets} & \textbf{LR} & \textbf{Dropout} & \textbf{WD} & \textbf{Unfreeze Epoch} \\
\midrule
1 & ESC-50, EmoDB & $1 \times 10^{-4}$ & 0.5 & 0.01 & 0 \\
2 & Italian PD & $5 \times 10^{-5}$ & 0.5 & 0.01 & 0 \\
3 & Pitt, PhysioNet & $1 \times 10^{-5}$ & 0.7 & 0.1 & 10 \\
\bottomrule
\end{tabular}
\end{table}

The alternative considered was a single universal configuration. This was rejected based on empirical evidence: the regularisation pendulum (underfitting on ESC-50 with clinical parameters, overfitting on Pitt with standard parameters) demonstrated that dataset-specific tuning is necessary when corpora span fundamentally different signal-to-noise ratios and sample sizes.

\section{Statistical Testing}

All primary comparisons use the Wilcoxon signed-rank test on 5-fold paired observations. This non-parametric test was chosen over the paired $t$-test because normality cannot be reliably verified with $N\!=\!5$ fold observations. The alternative considered, bootstrapped confidence intervals, was rejected as unnecessarily complex for the straightforward paired comparisons in this evaluation. Results are reported as mean $\pm$ standard deviation with $N$ explicitly stated.

\section{Reproducibility Statement}

The following measures ensure full reproducibility of all reported results.

\textbf{Fixed Seeds.} All experiments use seeds 42, 123, and 999. Each seed is set globally for Python \texttt{random}, NumPy, PyTorch CPU, and CUDA before every run via the \texttt{set\_seed()} function.

\textbf{Dataset Hashing.} Clinical datasets (Italian PD, Pitt Corpus) use SHA-256 cryptographic immutability tokens embedded in spectrogram filenames (e.g., \texttt{157\_2c87d4e1a0b5\_015\_orig.png}), linking every spectrogram to its source recording. Split manifests are locked as \texttt{splits\_manifest\_*.json} SHA-256 checksums.

\textbf{Pipeline Determinism.} \texttt{train\_unified.py} uses static ImageNet normalisation parameters ($\mu = [0.485, 0.456, 0.406]$, $\sigma = [0.229, 0.224, 0.225]$); no training-set statistics are computed at runtime, eliminating normalisation leakage.

\textbf{Leakage Guard.} An inline assertion in \texttt{train\_unified.py} verifies zero subject overlap between train and validation sets before every training run. A \texttt{RuntimeError} is raised and execution halts if any overlap is detected.

\textbf{Code Availability.} The full codebase is available in the project GitLab repository, including all split CSVs, training scripts, Kaggle notebooks, and result artefacts.

Few published studies on the Italian PD dataset report equivalent reproducibility guarantees.

% ═══════════════════════════════════════════════════════════════
% CHAPTER 5 — SYSTEM DESIGN AND ARCHITECTURE
% ═══════════════════════════════════════════════════════════════
\chapter{System Design and Architecture}

\section{Architectural Evolution}

The final \dualcnn{} architecture was not designed in isolation; it emerged from a six-phase experimental progression in which each design decision was motivated by empirical evidence from the preceding phase. Table~\ref{tab:phases} summarises this progression.

\begin{table}[H]
\centering
\caption{Architectural Evolution Across Six Phases}
\label{tab:phases}
\small
\begin{tabularx}{\textwidth}{c l X X}
\toprule
\textbf{Phase} & \textbf{Period} & \textbf{Question Asked} & \textbf{Key Finding} \\
\midrule
Term~1 & Oct--Dec 2025 & Which CNN baseline for spectrogram classification? & ResNet50 best; CBAM unstable ($\pm 19\%$); SE marginal \\
1 & Jan 2026 & Do baselines transfer to clinical audio? & EmoDB saturated (100\%); Pitt biased (66\%, C-Recall $<$46\%) \\
2 & Feb 2026 & K-fold baselines with 3 architectures $\times$ 5 datasets? & HybridNet wins non-clinical; ResNet50 wins Italian PD; Fold~3 anomaly confirmed \\
3--4 & Feb--Mar 2026 & Which attention mechanism best? Strategy A$\to$B pivot & CA best single module; AG adds 8M params for no gain; Fold~3 amplified by isotropic attention \\
5 & Mar 2026 & Can SpecAugment + BiLSTM + dual branch reach SOTA? & Stage A: 0.960; Stage B (3 seeds): 0.964 $\pm$ 0.051 \\
6 & Mar 2026 & Final SOTA push with FP-SA and label smoothing & 0.964 $\pm$ 0.048; strongest verifiable result on Italian PD under certified evaluation \\
\bottomrule
\end{tabularx}
\end{table}

\subsection{Term~1: Baseline Establishment and CBAM Instability}

The interim project evaluated Baseline CNN (trained from scratch), AlexNet, ResNet50, ResNet50+SE, and ResNet50+CBAM on ESC-50 and EmoDB. ResNet50 achieved the strongest results (ESC-50: 89.9\% $\pm$ 0.9\%; EmoDB: 95.7\% $\pm$ 3.2\%), while CBAM exhibited catastrophic variance (ESC-50: 73.4\% $\pm$ 10.9\%; EmoDB: 80.2\% $\pm$ 19.3\%). The CBAM instability was attributed to the sequential channel-then-spatial attention pipeline amplifying noise when applied to spectrogram inputs with relatively uniform spatial structure. This finding motivated the pivot away from CBAM toward Coordinate Attention in Term~2.

\subsection{Phase~1: Clinical Domain Extension}

Gold anchor baselines (ResNet50, 3-seed mean, 80/20 split) were established on five datasets. EmoDB saturation (100\% accuracy) confirmed that the task was too easy to differentiate architectures. The Pitt Corpus exhibited severe class bias (66.19\% accuracy, Control Recall $<$46\%), motivating the weighted loss bias-correction protocol. Italian PD achieved 97.52\% accuracy (F1\,=\,0.937) on the leakage-prone 80/20 split, a number that would need to be re-evaluated under speaker-independent K-fold.

\subsection{Phase~2: K-Fold Baseline Matrix}

The transition from 80/20 to speaker-independent 5-fold CV was the single most important methodological upgrade. 75 training runs (5 datasets $\times$ 3 models $\times$ 5 folds) revealed: HybridNet dominance on non-clinical datasets (ESC-50: 0.923, EmoDB: 0.976); ResNet50 superiority on Italian PD (0.916 $\pm$ 0.046); and the first confirmed observation of the Fold~3 anomaly, where all three models dropped to F1 $\approx$ 0.80--0.85 on Fold~3 of Italian PD, while maintaining $>$0.95 on Folds 0--2.

\subsection{Phase~3--4: Attention Ablation and the Strategy A$\to$B Pivot}

The initial attention ablation (Strategy~A) changed four variables simultaneously relative to the Phase~2 baseline: backbone unfreeze epoch (0$\to$7), learning rate ($10^{-4}\to 5\times 10^{-5}$), weighted loss (off$\to$on), and attention module (none$\to$CA/AG). This confounded design was identified and corrected: Strategy~B matched Phase~2 exactly with attention as the sole experimental variable. The confounded variables are shown in Table~\ref{tab:strategy_pivot}.

\begin{table}[H]
\centering
\caption{Strategy A$\to$B Pivot: Identifying Confounded Variables}
\label{tab:strategy_pivot}
\small
\begin{tabular}{l c c c}
\toprule
\textbf{Setting} & \textbf{Phase~2 Baseline} & \textbf{Strategy A} & \textbf{Strategy B} \\
\midrule
Backbone trainable from & Epoch 0 & Epoch 7 & \textbf{Epoch 0} \\
Learning rate & $1 \times 10^{-4}$ & $5 \times 10^{-5}$ & $\mathbf{1 \times 10^{-4}}$ \\
Weighted loss & No & Yes & \textbf{No} \\
Attention module & None & CA/Gate & \textbf{CA/Gate} \\
\bottomrule
\end{tabular}
\end{table}

Strategy~B results on Italian PD showed CA improving median F1 from 0.950 to 0.962 ($+1.2\%$) while the Attention Gate reduced mean F1 by 1.1\%. A critical observation was the architectural bug: the initial implementation injected 16 attention modules (one per residual block) into ResNet50, causing model collapse. The fix, a single module at the end of \texttt{layer4} only, was essential for stable training.

\subsection{Phase~5--6: SOTA Push}

SpecAugment, BiLSTM temporal modelling, and dual-branch fusion were introduced incrementally. The ablation table (Table~\ref{tab:ablation_results}) in Chapter~7 documents the contribution of each component.

\section{The \texttt{dual\_cnn\_sa\_lstm} Architecture}

The final architecture processes a $224 \times 224 \times 3$ log-Mel spectrogram through the following forward pass:

\begin{enumerate}
    \item \textbf{Dual-Branch Feature Extraction.}
    \begin{itemize}[nosep]
        \item \textbf{Branch 1:} ResNet50 (ImageNet-pretrained, fine-tuned from \texttt{layer4}) produces a $(B, 2048, 7, 7)$ feature map.
        \item \textbf{Branch 2:} EfficientNetV2-S (ImageNet-pretrained, fine-tuned) produces a $(B, 1280, 7, 7)$ feature map after adaptive average pooling to $7 \times 7$.
    \end{itemize}

    \item \textbf{Channel Concatenation and Projection.} Features are concatenated along the channel dimension: $(B, 3328, 7, 7)$. A $1 \times 1$ convolution with batch normalisation projects to $(B, 512, 7, 7)$, reducing dimensionality while fusing complementary feature representations.

    \item \textbf{FrequencyPriorSelfAttention (FP-SA).} Multi-head self-attention with 8 heads and reduction factor 4 operates on the $49$-token sequence ($7 \times 7$ spatial grid), with the learnable frequency-axis prior bias injected into the key computation. The output retains the $(B, 512, 7, 7)$ spatial structure via a residual connection.

    \item \textbf{Sequence Reshaping.} The spatial grid is flattened and permuted to $(B, 49, 512)$, a 49-step sequence where each step corresponds to a spatial location in the feature map.

    \item \textbf{BiLSTM Temporal Head.} A 2-layer bidirectional LSTM with hidden size 512 processes the sequence, producing $(B, 49, 1024)$. Mean pooling over the sequence dimension yields $(B, 1024)$.

    \item \textbf{Classifier.} Dropout (0.5) followed by a linear layer produces the output logits.
\end{enumerate}

The choice of EfficientNetV2-S as the second branch (rather than VGG16, MobileNetV2, or a second ResNet) was motivated by three considerations: (1)~compound scaling captures fine-grained frequency details that complement ResNet50's structural features; (2)~MobileNetV2 was empirically weaker on clinical datasets in Phase~2; (3)~a heterogeneous dual branch provides genuinely complementary representations, whereas two ResNets would produce highly correlated features.

\begin{figure}[htbp]
    \centering
    \includegraphics[width=\linewidth]{figures/fig_architecture.png}
    \caption{Architecture of the \dualcnn{} model showing the dual-branch feature extraction, FrequencyPriorSelfAttention (FP-SA), Coordinate Attention, and BiLSTM temporal head.}
    \label{fig:architecture}
\end{figure}

\section{HybridNet: $\alpha$-Gate Fusion}

The HybridNet architecture, which dominated non-clinical datasets in Phase~2, combines ResNet50 (2048-dim) and MobileNetV2 (1280-dim) features through a learnable sigmoid gate. Batch normalisation after concatenation prevents the higher-dimensional ResNet branch from dominating the fused representation. The gate vector $\alpha \in \mathbb{R}^{3328}$ is initialised to 0.5 and learned during training. This architecture served as the strongest non-clinical baseline and informed the dual-branch design of \dualcnn{}.

\section{Software Engineering Design}

\subsection{Model Factory Pattern}

\texttt{model\_builder.py} implements a factory pattern that constructs any backbone-attention combination from a single string argument: \texttt{resnet50\_ca}, \texttt{hybrid\_gate}, \texttt{dual\_cnn\_sa\_lstm}, etc. The factory parses compound attention types (e.g., \texttt{ca\_ag} for mixed CA+AG injection) and routes to the appropriate constructor. This design enables the 75+ experimental runs to be orchestrated via CLI arguments without code modification.

\subsection{Dataset-Agnostic Training Engine}

\texttt{train\_unified.py} provides a single entry point for all datasets. The \texttt{DS\_CONFIG} dictionary maps dataset names to their fold-indexed CSV paths, number of classes, and label mappings. The \texttt{--fold} argument routes to fold-specific CSVs; omitting it falls back to legacy single-split mode for backward compatibility. This design eliminated the need for dataset-specific training scripts and ensured identical training logic across all experiments.

\subsection{GitLab CI/CD Pipeline}

The \texttt{.gitlab-ci.yml} defines four stages: \texttt{setup} (dependency installation), \texttt{test} (Ruff linting, Black formatting, mypy type checking, pytest), \texttt{audit} (data leakage verification via \texttt{verify\_leakage.py}), and \texttt{package} (smoke test: one-epoch training on EmoDB to verify end-to-end pipeline integrity). This continuous integration pipeline was maintained throughout Term~2 development.

% ═══════════════════════════════════════════════════════════════
% CHAPTER 6 — IMPLEMENTATION
% ═══════════════════════════════════════════════════════════════
\chapter{Implementation}

\section{Directory Structure}

The codebase is organised into four principal directories:

\begin{verbatim}
product/
├── audio_preprocessing/
│   ├── data/                      # Raw audio (excluded from repo)
│   ├── src/                       # Spectrogram generation scripts
│   │   ├── audio_utils.py         # I/O, log-Mel, augmentation
│   │   ├── generate_spectrograms.py          # ESC-50
│   │   ├── generate_spectrograms_emodb.py    # EmoDB
│   │   ├── generate_spectrograms_italian.py  # Italian PD
│   │   ├── generate_spectrograms_pcgita.py   # PC-GITA
│   │   └── generate_spectrograms_pitt.py     # Pitt Corpus
│   └── outputs/                   # Generated spectrograms
├── models/
│   ├── dual_cnn_sa_lstm.py        # Novel architecture
│   ├── freq_prior_attention.py    # FP-SA module
│   ├── coordinate_attention.py    # CA module
│   ├── model_builder.py           # Factory pattern
│   ├── hybrid_net.py              # HybridNet
│   ├── resnet_bilstm.py           # ResNet+BiLSTM
│   └── [se_block, cbam, attention_gate, ...]
├── training/
│   └── train_unified.py           # Training engine
├── artifacts/
│   ├── splits/                    # All fold CSVs
│   └── runs/                      # Results, checkpoints
└── notebooks/                     # Kaggle notebooks
\end{verbatim}

\section{The Unified Training Engine}

\texttt{train\_unified.py} accepts the following CLI arguments: \texttt{--dataset} (choice of six datasets), \texttt{--model\_type} (compound backbone\_attention string), \texttt{--fold} (0--4 for K-fold mode), \texttt{--seed}, \texttt{--epochs}, \texttt{--lr}, \texttt{--dropout}, \texttt{--weight\_decay}, \texttt{--unfreeze\_at}, \texttt{--weighted\_loss}, \texttt{--spec\_augment}, \texttt{--label\_smoothing}, and optional CSV path overrides. The engine handles label mapping, class weight computation, transfer learning freeze/unfreeze scheduling, TensorBoard logging, confusion matrix generation, and JSON summary serialisation.

The model type string is parsed at a single underscore boundary: \texttt{resnet50\_ca} splits into backbone \texttt{resnet50} and attention \texttt{ca}. Compound types like \texttt{dual\_cnn\_sa\_lstm} are intercepted by the extended factory before the generic parser runs.

\section{Kaggle Notebook Pipeline}

GPU-accelerated training was executed on Kaggle (Tesla T4/P100). Notebooks were generated programmatically from templates, with each notebook targeting a specific dataset-model-seed-fold combination. The notebook pipeline patches \texttt{train\_unified.py} and \texttt{model\_builder.py} at runtime to fix Kaggle-specific path resolution issues.

The Phase~5 Kaggle leakage bug warrants detailed documentation. The root cause was Kaggle dataset versioning: when the dataset was re-uploaded with corrected fold CSVs, Kaggle retained stale versions from prior uploads. Folds~0 and~1 resolved to corrupted CSVs from an earlier version, producing artificially perfect F1\,=\,1.0, while Folds~2--4 loaded the correct CSVs and produced realistic scores (0.85--0.93). The bimodal distribution was the diagnostic signal. The fix comprised three elements: (1)~hardcoding the dataset slug \texttt{phase6-italian-pd-clean} to bypass Kaggle's versioned path resolution; (2)~replacing the generic path-fix cell with an inline leakage guard; (3)~repacking and re-uploading the dataset as a fresh entity.

\section{SHA-256 Cryptographic Immutability}

Clinical datasets use SHA-256 hashes embedded in spectrogram filenames: \texttt{157\_2c87d4e1a0b5\_015\_orig.png}. This provides cryptographic traceability from any spectrogram back to its source recording, enabling post-hoc audit of any result. Split manifests are locked as \texttt{splits\_manifest\_*.json} checksums, preventing silent modification of fold assignments.

\section{CI/CD Pipeline}

The GitLab pipeline (\texttt{.gitlab-ci.yml}) executes four stages on every push: \texttt{setup} (Python 3.12, PyTorch CPU, project dependencies); \texttt{lint} (Ruff, Black, mypy); \texttt{audit} (data leakage verification); \texttt{package} (one-epoch smoke test on EmoDB). The smoke test produces JSON and PNG artefacts that are retained for one week, providing continuous verification that the training pipeline remains functional. Fourteen Ruff linting errors were identified and fixed during initial CI setup.

\section{TensorBoard Logging and Artefact Generation}

Every training run logs scalar metrics (train/val loss, accuracy, macro F1, control recall, AUC) to TensorBoard. Upon completion, each run saves: \texttt{summary.json} (best and final metrics with full configuration), \texttt{best\_classification\_report.json} (per-class precision, recall, F1), \texttt{confusion\_matrix.png}, and \texttt{best\_model.pt} (state dictionary of the best-F1 checkpoint). This standardised artefact structure enabled automated aggregation across 75+ runs.

% ═══════════════════════════════════════════════════════════════
% CHAPTER 7 — RESULTS AND ANALYSIS
% ═══════════════════════════════════════════════════════════════
\chapter{Results and Analysis}

This chapter presents results in chronological experimental order, with theoretical analysis accompanying every numerical finding. All results use exact values from experimental logs.

\section{Term~1 Revisited: Baseline Establishment}

Term~1 established CNN baselines on ESC-50 and EmoDB under 80/20 splitting with three seeds. Table~\ref{tab:term1} summarises the key findings.

\begin{table}[H]
\centering
\caption{Term~1 Baseline Results (80/20 split, 3 seeds)}
\label{tab:term1}
\begin{tabular}{l c c c c}
\toprule
\textbf{Model} & \textbf{ESC-50 Acc} & \textbf{ESC-50 F1} & \textbf{EmoDB Acc} & \textbf{EmoDB F1} \\
\midrule
Baseline CNN & 44.4\% & 0.424 & --- & --- \\
AlexNet & 74.5\% $\pm$ 4.3\% & 0.738 & --- & --- \\
ResNet50 & 89.9\% $\pm$ 0.9\% & 0.898 & 95.7\% $\pm$ 3.2\% & 0.946 \\
ResNet50 + SE & 88.8\% $\pm$ 2.4\% & 0.887 & 95.8\% $\pm$ 0.9\% & 0.949 \\
ResNet50 + CBAM & 73.4\% $\pm$ 10.9\% & --- & 80.2\% $\pm$ 19.3\% & --- \\
\bottomrule
\end{tabular}
\end{table}

The critical finding was CBAM's catastrophic instability: variance of $\pm 10.9\%$ on ESC-50 and $\pm 19.3\%$ on EmoDB across three seeds. This instability arises because CBAM's sequential channel-then-spatial attention pipeline amplifies noise in the spatial branch when operating on spectrogram inputs whose spatial structure (frequency $\times$ time) is more uniform than the object-centric scenes ImageNet attention modules were designed for. The spatial attention map, computed via max and average pooling across channels, produces near-uniform weights on spectrograms, causing the sigmoid output to oscillate near 0.5 and destabilise gradient flow. The instability motivated the pivot to Coordinate Attention, which avoids the spatial branch entirely in favour of factorised directional pooling.

\section{Phase~1: Clinical Domain Extension}

Gold anchor baselines (ResNet50, 3-seed mean, 80/20 split) were established across five datasets. Table~\ref{tab:phase1} presents these anchors.

\begin{table}[H]
\centering
\caption{Phase~1 Gold Anchor Baselines (ResNet50, 3-seed mean, 80/20 split)}
\label{tab:phase1}
\begin{tabular}{l c c c}
\toprule
\textbf{Dataset} & \textbf{Accuracy} & \textbf{Macro F1} & \textbf{AUC} \\
\midrule
Italian PD & 97.52\% & 0.937 & 0.992 \\
PhysioNet & 93.13\% & 0.884 & 0.955 \\
ESC-50 & 74.67\% & 0.731 & 0.987 \\
EmoDB & 100.00\% & 1.000 & 1.000 \\
Pitt Corpus & 66.19\% & 0.577 & 0.621 \\
\bottomrule
\end{tabular}
\end{table}

EmoDB saturation at 100\% confirmed the task is too easy for architecture differentiation; the 10-speaker dataset with strong emotional prosody provides insufficient challenge for ImageNet-pretrained ResNet50. This observation justified the pivot to clinical speech where the signal-to-noise ratio between pathological and healthy speech is substantially lower.

The Pitt Corpus baseline (66.19\% accuracy, AUC\,=\,0.621) initially appeared to indicate model failure. Investigation revealed severe class bias: the model was predicting the majority class (dementia) for nearly all samples, achieving inflated accuracy on the imbalanced dataset while failing to identify 54\% of healthy controls. The bias-correction protocol (weighted cross-entropy loss with class weights inversely proportional to class frequency) was implemented to address this, and Macro F1 was adopted as the primary metric specifically because it is robust to class imbalance.

\section{Phase~2: K-Fold Baseline Matrix}

The 75-run K-fold matrix ($5 \text{ datasets} \times 3 \text{ models} \times 5 \text{ folds}$) constitutes the most comprehensive baselines in this thesis. Table~\ref{tab:phase2} presents the aggregated results.

\begin{table}[H]
\centering
\caption{Phase~2 K-Fold Baseline Matrix (Macro F1, Mean $\pm$ Std, 5-fold)}
\label{tab:phase2}
\begin{tabular}{l c c c c}
\toprule
\textbf{Dataset} & \textbf{ResNet50} & \textbf{MobileNetV2} & \textbf{HybridNet} & \textbf{Winner} \\
\midrule
EmoDB & 0.942 $\pm$ 0.017 & 0.856 $\pm$ 0.035 & \textbf{0.976 $\pm$ 0.006} & HybridNet \\
ESC-50 & 0.835 $\pm$ 0.014 & 0.826 $\pm$ 0.027 & \textbf{0.923 $\pm$ 0.006} & HybridNet \\
PhysioNet & 0.889 $\pm$ 0.013 & 0.884 $\pm$ 0.008 & \textbf{0.891 $\pm$ 0.011} & HybridNet \\
Italian PD & \textbf{0.916 $\pm$ 0.046} & 0.911 $\pm$ 0.063 & 0.908 $\pm$ 0.061 & ResNet50 \\
\bottomrule
\end{tabular}
\end{table}

HybridNet's dominance on non-clinical datasets (ESC-50: $+8.8\%$ over ResNet50; EmoDB: $+3.4\%$) is explained by its $\alpha$-gate fusion mechanism: the learnable gate adaptively combines ResNet50's high-level semantic features (2048-dim) with MobileNetV2's efficient spatial features (1280-dim), providing complementary representations for tasks with diverse spectral patterns. The batch normalisation scale alignment prevents the higher-dimensional ResNet branch from dominating, allowing the fusion to draw on both streams effectively.

However, on Italian PD, ResNet50 alone outperformed HybridNet (0.916 vs 0.908). PD detection from spectrograms is fundamentally a fine-grained texture classification task: the pathological signal is a subtle shift in formant structure, not a change in scene composition. ResNet50's deep residual features capture these texture variations more effectively than MobileNetV2's depthwise separable convolutions, which are optimised for efficient spatial feature extraction at the expense of channel-wise expressiveness. The fusion overhead of HybridNet (3328 combined dimensions, 512-unit MLP, 0.7 dropout) introduces unnecessary capacity for this binary texture discrimination task.

\subsection{Fold~3 Anomaly: First Confirmed Observation}
\label{subsec:fold3}

Phase~2 produced an unexpected finding. All three architectures (ResNet50, MobileNetV2, and HybridNet) show consistent, substantial performance degradation on Fold~3 of the Italian PD dataset. The pattern is architectural-agnostic: three models with fundamentally different designs all degrade on the same fold, which rules out any model-specific weakness and points to a data-level cause. Table~\ref{tab:fold3_phase2} shows the fold-level breakdown.

\begin{table}[H]
\centering
\caption{Phase~2: Italian PD Fold-Level F1 Scores}
\label{tab:fold3_phase2}
\begin{tabular}{c c c c}
\toprule
\textbf{Fold} & \textbf{ResNet50} & \textbf{MobileNetV2} & \textbf{HybridNet} \\
\midrule
0 & 0.950 & 0.983 & 0.978 \\
1 & 0.981 & 0.994 & 0.975 \\
2 & 0.957 & 0.965 & 0.957 \\
\textbf{3} & \textbf{0.851} & \textbf{0.797} & \textbf{0.797} \\
4 & 0.892 & 0.910 & 0.904 \\
\bottomrule
\end{tabular}
\end{table}

\begin{figure}[htbp]
    \centering
    \includegraphics[width=\linewidth]{figures/fig_fold3_anomaly.png}
    \caption{Phase~2 fold-level F1 scores on the Italian PD dataset across three architectures. The consistent drop at Fold~3 is visible across all models, with the degradation substantially larger than any other fold variation. This consistency across architecturally distinct models indicates a data-level rather than model-level phenomenon.}
    \label{fig:fold3_anomaly}
\end{figure}

The Fold~3 drop is sharp and consistent across all three models, not a gradual variation. The finding was subsequently confirmed across all three random seeds (42, 123, 999) in Phase~5 experiments (see Table~\ref{tab:fold3_seed_consistency} in Appendix~C). Seed invariance rules out initialisation sensitivity as an explanation: the difficulty is a structural property of the Fold~3 validation set composition.

Zero data leakage was confirmed for all Phase~2 splits. Investigation revealed that Fold~3's validation set contains approximately 21 healthy control subjects whose speech spectrograms consistently resemble PD patterns: HC recall drops to approximately 75\% while PD recall remains 1.000. The model identifies every PD patient correctly; the confusion is entirely localised to a specific subgroup of HC subjects who produce PD-like spectrograms despite carrying no clinical PD diagnosis at the time of recording. Single-split evaluation, which is standard in the literature, would have made this fold-level instability entirely invisible. To the best of the author's knowledge, this borderline HC cluster has not been previously identified or characterised in published work on this dataset.

\section{Phase~3--4: Attention Ablation}

Strategy~B isolates attention as the sole experimental variable. Table~\ref{tab:ablation_attention} presents the controlled comparison on Italian PD.

\begin{table}[H]
\centering
\caption{Phase~4 Attention Ablation, Strategy~B (Italian PD, 1 seed)}
\label{tab:ablation_attention}
\begin{tabular}{c c c c c c}
\toprule
\textbf{Fold} & \textbf{ResNet50} & \textbf{+CA} & $\boldsymbol{\Delta}$ & \textbf{+Gate} & $\boldsymbol{\Delta}$ \\
\midrule
0 & 0.950 & 0.972 & $+2.2\%$ & 0.956 & $+0.6\%$ \\
1 & 0.981 & 0.962 & $-1.9\%$ & 0.975 & $-0.6\%$ \\
2 & 0.957 & 0.965 & $+0.9\%$ & 0.974 & $+1.8\%$ \\
3 & 0.851 & 0.818 & $-3.4\%$ & 0.811 & $-4.1\%$ \\
4 & 0.892 & 0.882 & $-1.0\%$ & 0.863 & $-2.9\%$ \\
\midrule
\textbf{Mean} & \textbf{0.926} & \textbf{0.920} & $-0.7\%$ & \textbf{0.916} & $-1.1\%$ \\
\textbf{Median} & \textbf{0.950} & \textbf{0.962} & $+1.2\%$ & \textbf{0.956} & $+0.6\%$ \\
\bottomrule
\end{tabular}
\end{table}

\begin{figure}[htbp]
    \centering
    \includegraphics[width=0.9\linewidth]{figures/fig_attention_ablation.png}
    \caption{Phase 4 attention ablation (Strategy B) showing per-fold F1 scores on the Italian PD dataset. Coordinate Attention (CA) provides the most stable improvement across folds compared to standard ResNet50 and Attention Gate (Gate).}
    \label{fig:attention_ablation}
\end{figure}

The divergence between mean and median effects is directly attributable to Fold~3. Attention improves performance on the ``easy'' folds (0, 2) where the data partition provides clean HC/PD separation, but \emph{amplifies} the degradation on Fold~3. Isotropic attention mechanisms redistribute representational capacity uniformly across all spatial positions. On Fold~3, where the HC borderline cluster presents genuinely ambiguous spectral features, this redistribution amplifies the confusion rather than resolving it. This was a key motivation for FP-SA: by injecting a frequency-axis prior, the attention mechanism is directed toward the frequency bands where PD pathology manifests, rather than spreading attention uniformly.

The Attention Gate adds approximately 8M parameters (23.4M total vs 15.4M for CA) for worse performance than CA on both mean and median F1. The additional capacity provides no benefit because the spatial suppression mask is poorly suited to spectrogram inputs: the ``irrelevant regions'' that the gate learns to suppress in medical images (background tissue) have no clear analogue in spectrograms where the entire time-frequency plane carries potential diagnostic information.

\section{Phase~4: Cross-Lingual Evaluation on PC-GITA}

To assess cross-lingual generalisation, the proposed architecture was evaluated on the Colombian Spanish PC-GITA diadochokinetic (DDK) corpus. The \dualcnn{} model achieves an overall F1 score of $0.780 \pm 0.056$ under $N\!=\!15$ speaker-independent cross-validation. 

Unlike the Italian PD dataset which features read speech, PC-GITA DDK tasks stress the articulatory subsystem through rapid syllable repetition. The model achieves an AUC of 0.822, with a severe drop in healthy control recall (0.517) compared to PD recall. The fold-level F1 breakdown (Fold~0: 0.754, Fold~1: 0.800, Fold~2: 0.830, Fold~3: 0.691, Fold~4: 0.824) demonstrates significant variance across the diverse Colombian Spanish demographic. This cross-lingual gap confirms that while deep spectrogram features capture underlying pathological markers, language-specific phonological constraints and task differences heavily influence generalisation.

\section{Phase~5--6: SOTA Push and Ablation}

\subsection{Incremental Ablation on Italian PD}

Table~\ref{tab:ablation_results} presents the controlled ablation of the final architecture on Italian PD, adding one component at a time.

\begin{table}[H]
\centering
\caption{Italian PD Ablation (seed 42, 5-fold)}
\label{tab:ablation_results}
\begin{tabular}{l c c}
\toprule
\textbf{Model} & \textbf{Mean F1 $\pm$ Std} & \textbf{$\Delta$ vs previous} \\
\midrule
ResNet50 baseline & 0.950 $\pm$ 0.049 & --- \\
+ Coordinate Attention & 0.957 $\pm$ 0.041 & $+0.7\%$ \\
+ BiLSTM temporal head & 0.958 $\pm$ 0.058 & $+0.1\%$ \\
+ EfficientNetV2-S branch & 0.964 $\pm$ 0.048 & $+0.6\%$ \\
Full \dualcnn{} & 0.964 $\pm$ 0.048 & confirmed \\
\bottomrule
\end{tabular}
\end{table}

\begin{figure}[htbp]
    \centering
    \includegraphics[width=0.8\linewidth]{figures/fig_ablation.png}
    \caption{Ablation study on the Italian PD dataset (5-fold mean F1, seed 42). The incremental addition of Coordinate Attention, BiLSTM, and the EfficientNetV2-S branch improves performance towards the published state of the art.}
    \label{fig:ablation}
\end{figure}

Each component contributes measurably. CA provides $+0.7\%$ by enabling direction-aware attention on the frequency and time axes independently, since PD formant shifts are primarily frequency-axis phenomena, and CA can learn to weight this axis preferentially. The BiLSTM adds $+0.1\%$: a modest but consistent gain from temporal modelling across the preserved time columns. The EfficientNetV2-S branch adds $+0.6\%$ by providing complementary fine-grained features via compound scaling that ResNet50's fixed-scale residual blocks cannot capture.

The variance also drops from 0.049 (baseline) to 0.041 with CA, reflecting more consistent fold-to-fold behaviour. CA gives the network a learnable mechanism to weight the most diagnostically informative spatial regions, which reduces sensitivity to fold-specific data variation. The variance increase with BiLSTM (0.058) reflects the LSTM's sensitivity to sequence-level patterns that vary more across folds than spatial features.

\subsection{Stage~A: Single-Seed Results}

Table~\ref{tab:stage_a} shows the single-seed SOTA push results on Italian PD with SpecAugment.

\begin{table}[H]
\centering
\caption{Phase~5 Stage~A (Italian PD, 1 seed, 5-fold, SpecAugment)}
\label{tab:stage_a}
\begin{tabular}{l c c c}
\toprule
\textbf{Model} & \textbf{Mean F1 $\pm$ Std} & \textbf{AUC} & \textbf{Control Recall} \\
\midrule
ResNet50 + SpecAug & 0.937 $\pm$ 0.053 & 0.980 $\pm$ 0.027 & 0.884 $\pm$ 0.117 \\
ResNet50 + CA + SpecAug & 0.948 $\pm$ 0.055 & 0.985 $\pm$ 0.019 & 0.893 $\pm$ 0.121 \\
ResNet50 + CA + BiLSTM & 0.951 $\pm$ 0.059 & 0.988 $\pm$ 0.017 & 0.926 $\pm$ 0.150 \\
\texttt{dual\_cnn\_lstm} + SpecAug & \textbf{0.960 $\pm$ 0.043} & 0.988 $\pm$ 0.016 & \textbf{0.928 $\pm$ 0.091} \\
\bottomrule
\end{tabular}
\end{table}

The dual-branch model achieves both the highest mean F1 and the lowest variance, which is consistent with the complementary feature representations from ResNet50 and EfficientNetV2-S producing a more stable basis for clinical classification than either branch alone.

\subsection{Stage~B: Final Multi-Seed Result}

The final result on Italian PD with 3 seeds $\times$ 5 folds ($N\!=\!15$) and label smoothing ($\varepsilon\!=\!0.05$):

\begin{equation}
    \text{F1}_{\texttt{dual\_cnn\_sa\_lstm}} = 0.964 \pm 0.048 \quad (N = 15)
\end{equation}

Label smoothing was applied at $\varepsilon\!=\!0.05$ to soften hard label boundaries, particularly for the ambiguous Fold~3 HC subjects. The alternative considered was $\varepsilon\!=\!0.1$, which was rejected as too aggressive for a binary task where one class (PD) shows near-perfect separability on most folds.

\subsection{Statistical Significance}
\label{subsec:wilcoxon}

To validate the improvement of the final \dualcnn{} over the ResNet50 baseline, a paired Wilcoxon signed-rank test was conducted. Both models were evaluated under identical conditions: speaker-independent 5-fold cross-validation on the Italian PD dataset across all three random seeds (42, 123, 999), yielding $N\!=\!15$ paired observations per model (one per seed-fold combination). Each pair $(x_i^{\text{dual}},\, x_i^{\text{base}})$ corresponds to the same seed and fold, controlling for both data partitioning and initialisation variability.

Table~\ref{tab:wilcoxon_pairs} presents the 15 paired fold-level F1 scores. The signed differences $d_i = x_i^{\text{dual}} - x_i^{\text{base}}$ are ranked by absolute value; the Wilcoxon statistic $W$ is the smaller of the sum of positive- and negative-signed ranks.

\begin{table}[H]
\centering
\caption{Wilcoxon Paired Observations: \dualcnn{} vs ResNet50 Baseline (Italian PD, $N\!=\!15$)}
\label{tab:wilcoxon_pairs}
\small
\begin{tabular}{c c c c c c}
\toprule
\textbf{Seed} & \textbf{Fold} & \textbf{ResNet50 F1} & \textbf{\dualcnn{} F1} & \textbf{$d_i$} & \textbf{Signed Rank} \\
\midrule
42  & 0 & 0.950 & 0.994 & $+$0.0441 & $+$9.5 \\
42  & 1 & 0.981 & 1.000 & $+$0.0189 & $+$3.0 \\
42  & 2 & 0.957 & 0.991 & $+$0.0347 & $+$6.0 \\
42  & 3 & 0.851 & 0.865 & $+$0.0133 & $+$1.0 \\
42  & 4 & 0.892 & 0.954 & $+$0.0621 & $+$15.0 \\
123 & 0 & 0.950 & 1.000 & $+$0.0496 & $+$13.0 \\
123 & 1 & 0.981 & 1.000 & $+$0.0189 & $+$3.0 \\
123 & 2 & 0.957 & 1.000 & $+$0.0434 & $+$7.5 \\
123 & 3 & 0.851 & 0.898 & $+$0.0470 & $+$12.0 \\
123 & 4 & 0.892 & 0.945 & $+$0.0535 & $+$14.0 \\
999 & 0 & 0.950 & 0.994 & $+$0.0441 & $+$9.5 \\
999 & 1 & 0.981 & 1.000 & $+$0.0189 & $+$3.0 \\
999 & 2 & 0.957 & 1.000 & $+$0.0434 & $+$7.5 \\
999 & 3 & 0.851 & 0.884 & $+$0.0329 & $+$5.0 \\
999 & 4 & 0.892 & 0.937 & $+$0.0447 & $+$11.0 \\
\midrule
\multicolumn{2}{l}{\textbf{Mean}} & $0.926 \pm 0.053$ & $\mathbf{0.964 \pm 0.048}$ & & \\
\bottomrule
\end{tabular}
\end{table}

All 15 signed differences are positive ($W^{+} = 120$, $W^{-} = 0$), giving a test statistic of $W = 0$ and an exact two-tailed $p = 6.1 \times 10^{-5}$ (one-tailed $p = 3.1 \times 10^{-5}$). The null hypothesis that the median difference is zero is rejected at $p < 0.0001$. The non-parametric test was chosen because normality of the difference distribution cannot be verified with only $N\!=\!5$ fold observations per seed. With every paired fold showing a positive difference and the smallest margin being $+0.013$ (Fold~3, seed~42), the improvement is not marginal: the \dualcnn{} architecture consistently outperforms the ResNet50 baseline across all three seeds and all five folds.

\section{Cross-Dataset Summary}

Table~\ref{tab:cross_dataset} presents the final cross-dataset performance.

\begin{table}[H]
\centering
\caption{Final Cross-Dataset Summary}
\label{tab:cross_dataset}
\begin{tabular}{l l c l}
\toprule
\textbf{Dataset} & \textbf{Best Model} & \textbf{Mean F1 $\pm$ Std} & \textbf{Notes} \\
\midrule
Italian PD & \dualcnn{} & 0.964 $\pm$ 0.048 & Primary clinical, $N\!=\!15$ \\
PC-GITA & \dualcnn{} & $0.780 \pm 0.056$ & Cross-lingual DDK, $N\!=\!15$ \\
ESC-50 & HybridNet & 0.923 $\pm$ 0.006 & Non-clinical, $N\!=\!15$ \\
EmoDB & HybridNet & 0.976 $\pm$ 0.006 & Non-clinical emotion, $N\!=\!15$ \\
\bottomrule
\end{tabular}
\end{table}

\begin{figure}[htbp]
    \centering
    \includegraphics[width=\linewidth]{figures/fig_cross_dataset.png}
    \caption{Cross-dataset performance comparison showing the improvement from the ResNet50 baseline to the final models. The \dualcnn{} architecture achieves competitive performance on Italian PD and establishes a benchmark on the PC-GITA cross-lingual task.}
    \label{fig:cross_dataset}
\end{figure}

The cross-lingual gap (0.964 $\to$ 0.780 on PC-GITA) quantifies the generalisation challenge. PC-GITA uses DDK tasks (syllable repetition) rather than read speech, and the corpus is in Colombian Spanish rather than Italian. The gap reflects the combined effect of: (1)~language-specific phonological constraints on PD manifestation; (2)~fundamentally different speech task types (read speech captures prosodic features; DDK captures articulatory motor control); (3)~the smaller PC-GITA corpus (600 WAVs vs 831 spectrograms) reducing the training data available per fold. Achieving 0.780 F1 on an entirely unseen language and task type suggests the learned spectrogram features transfer at least partially across languages.

\section{Comparison with Prior Work}

Table~\ref{tab:sota_comparison} compares this work against published results on the Italian PD dataset and related corpora. For each study, the evaluation protocol is examined across four criteria: speaker-independent splitting, formal leakage controls, multi-seed/multi-fold reporting, and direct comparability to the protocol used here. Where information was not disclosed in the source paper, the cell is marked ``not reported.''

\begin{table}[H]
\centering
\caption{Comparison with Prior Work on PD Speech Detection}
\label{tab:sota_comparison}
\small
\begin{tabularx}{\textwidth}{l l c c c c c}
\toprule
\textbf{Study} & \textbf{Dataset} & \textbf{Metric} & \textbf{Spkr} & \textbf{Leakage} & \textbf{Seeds/} & \textbf{Comp-} \\
 & & & \textbf{Ind.?} & \textbf{Ctrl?} & \textbf{Folds} & \textbf{arable?} \\
\midrule
Sakar et al.~\cite{sakar2013} & Own corpus & 75--85\% acc & NR & No & NR & No \\
Tsanas et al.~\cite{tsanas2012} & Proprietary & Higher & \textbf{No} & No & NR & No \\
Dimauro et al.~\cite{dimauro2019} & Italian PD & NR & NR & NR & NR & Uncertain \\
Aversano et al.~\cite{aversano2024} & Italian PD & F1\,=\,0.971 & NR & NR & NR & \textbf{No} \\
Orozco et al.~\cite{orozco2016} & Multi-lingual & Varies & NR & NR & NR & No \\
\midrule
\textbf{This work} & \textbf{Italian PD} & \textbf{F1\,=\,0.964\,$\pm$\,0.048} & \textbf{Yes} & \textbf{Yes} & \textbf{3$\times$5} & \textbf{Yes} \\
\textbf{This work} & \textbf{PC-GITA} & \textbf{F1\,=\,0.780\,$\pm$\,0.056} & \textbf{Yes} & \textbf{Yes} & \textbf{3$\times$5} & Ref. \\
\bottomrule
\multicolumn{7}{l}{\small NR = Not Reported in source paper.}
\end{tabularx}
\end{table}

The proposed system achieves F1\,=\,$0.964 \pm 0.048$ ($N\!=\!15$) on the Italian PD dataset under speaker-independent 5-fold cross-validation with certified zero leakage. The closest published figure on this dataset is Aversano et al.~\cite{aversano2024} (F1\,=\,0.971), but their paper does not specify whether speaker-independent splits were used, whether any leakage controls were applied, or how many seeds or folds produced the reported figure. Without this information, the two numbers cannot be meaningfully compared. A model evaluated without speaker-level separation can memorise speaker identity rather than pathological features, inflating performance in ways that do not generalise to unseen patients. The 0.007 gap is therefore not interpretable as a performance difference between the two systems: this work's result is produced under strictly documented and verifiable conditions, while the conditions behind the higher figure remain unknown. To the best of the author's knowledge, F1\,=\,0.964 is the strongest result on the Italian PD dataset reported under a fully reproducible, speaker-independent, zero-leakage evaluation protocol.

% ═══════════════════════════════════════════════════════════════
% CHAPTER 8 — DISCUSSION
% ═══════════════════════════════════════════════════════════════
\chapter{Discussion}

\section{Limitations}

Four specific limitations constrain the claims of this work.

First, the Italian PD corpus contains approximately 65 subjects. While the $N\!=\!15$ multi-seed evaluation provides statistical power for within-dataset conclusions, the small subject pool limits the external generalisability of absolute performance claims. Clinical deployment would require validation on substantially larger and more diverse populations.

Second, the Fold~3 borderline HC cluster has not been clinically characterised. The finding that approximately 21 HC subjects consistently produce PD-like spectrograms is robust (confirmed across all architectures and seeds), but the clinical interpretation (whether these subjects are pre-symptomatic PD, have other vocal disorders, or represent statistical outliers) cannot be determined without clinical follow-up data that is not available in the dataset.

Third, the FP-SA frequency prior is hand-designed: the functional form (inverse distance weighting on frequency bin indices) encodes a specific inductive bias that may not be optimal for every corpus. A data-driven learned prior could potentially adapt to corpus-specific spectral structures.

Fourth, no test-time augmentation or ensemble inference was implemented. The reported results reflect single-model inference with deterministic forward passes, leaving potential performance gains from inference-time techniques unexplored.

\section{Threats to Validity}

\subsection{Internal Validity}

Internal validity concerns whether experimental conditions are sufficiently controlled to draw causal conclusions. Both leakage bugs discovered during this project were caught and corrected before final results were generated. The inline RuntimeError guard provides an active ongoing defence against future leakage introduction. The Strategy~A$\to$B pivot explicitly corrected a confounded experimental design, ensuring that the ablation results in Table~\ref{tab:ablation_attention} reflect the isolated effect of attention modules.

One unresolved internal validity concern is the hand-designed FP-SA prior: while the frequency bias improves median F1, a fully controlled ablation between hand-designed and data-driven priors has not been conducted. The frequency prior represents a design bias that is not fully ablated. This is acknowledged as a limitation.

\subsection{External Validity}

External validity concerns whether results generalise beyond this study. The Italian PD corpus is limited to Italian speakers performing read speech tasks. PC-GITA covers Colombian Spanish DDK tasks, a different language and fundamentally different speech task. The cross-lingual gap (0.964 $\to$ 0.780) shows that generalisation across linguistic and demographic groups cannot be assumed. Both datasets are small by deep learning standards ($\sim$65 and 101 subjects respectively), further limiting external validity claims. Any deployment beyond the evaluated populations would require prospective validation.

\subsection{Construct Validity}

Construct validity concerns whether the metrics measure what is clinically relevant. Macro F1 is class-imbalance robust and clinically interpretable for binary detection tasks. However, it does not directly measure clinical utility. A model with F1\,=\,0.964 that systematically misclassifies a specific subgroup (as Fold~3 demonstrates) may be clinically unacceptable despite high aggregate performance. Control Recall is reported precisely because false positives carry significant clinical cost, since healthy patients receiving a PD diagnosis face unnecessary anxiety, medical interventions, and economic burden. The Fold~3 finding exposes a construct validity gap between aggregate metrics and subgroup-level clinical safety.

\subsection{Statistical Conclusion Validity}

$N\!=\!15$ per dataset (3 seeds $\times$ 5 folds) is sufficient for Wilcoxon signed-rank testing on 5-fold paired comparisons. Non-parametric testing was chosen because normality cannot be assumed with $N\!=\!5$ fold observations. Standard deviations reported reflect genuine fold-to-fold variance, not noise. The Fold~3 anomaly inflates Italian PD variance ($\text{std}\!=\!0.048$ vs $\text{std}\!=\!0.006$ for EmoDB). This is a data property reported transparently, not suppressed.

\section{Failure Case Analysis}

The Italian PD Fold~3 confusion matrix pattern is consistent across all architectures and all seeds:

\begin{itemize}[nosep]
    \item PD recall\,=\,1.000 (no Parkinson's patients missed)
    \item HC recall $\approx$ 0.75 (approximately 1 in 4 healthy controls misclassified as PD)
\end{itemize}

The pattern is consistent with the approximately 21 HC subjects in Fold~3 whose spectrograms consistently resemble PD patterns, rather than any architectural limitation. These subjects appear to represent a borderline subgroup, healthy individuals whose speech spectrograms exhibit PD-like spectral characteristics across all architectures and all random seeds, and this subgroup does not appear to have been previously identified in the literature.

The clinical implications are significant. If these HC subjects are pre-symptomatic PD (speech changes preceding motor diagnosis), the model is detecting a genuine pathological signal that the clinical labels have not yet captured. If they have other vocal disorders unrelated to PD, the model is correctly identifying spectral abnormality but incorrectly attributing it to PD. Either interpretation has important consequences for deployment.

To the best of our knowledge, no prior published work on this dataset reports fold-wise instability analysis or characterises this borderline HC cluster.

Grad-CAM visualisation of these misclassified HC spectrograms, identifying which spectral regions the model attends to when classifying these subjects as PD, would be a valuable extension that could disambiguate between the pre-symptomatic and alternative-disorder hypotheses. This is flagged as priority future work.

\section{Clinical Deployment Implications}

Responsible deployment of this system would require: prospective validation on unseen populations from multiple linguistic and demographic backgrounds; formal clinical trials with appropriate statistical powering; and regulatory approval (UKCA/CE marking for medical AI devices).

The Fold~3 finding has direct deployment implications: any deployment without explicit acknowledgement and safeguarding of the hard borderline cluster would be clinically irresponsible. A screening system that generates false positive PD diagnoses for a known subgroup of healthy individuals inflicts real harm, including anxiety, unnecessary follow-up procedures, and potential treatment of a non-existent condition.

The recommended deployment model is decision support rather than autonomous diagnosis: the system's output should be presented to a clinician as an additional data point alongside standard clinical assessment, not as a standalone diagnostic tool.

\section{Future Work}

Five directions for future research emerge from this work:

\begin{enumerate}
    \item \textbf{Learned Frequency Priors.} Replace the hand-designed FP-SA bias with a data-driven prior learned end-to-end during training, potentially using a small auxiliary network that predicts the optimal frequency weighting from the input spectrogram.

    \item \textbf{Grad-CAM Attention Visualisation.} Apply Grad-CAM to the Fold~3 misclassified HC spectrograms to identify which spectral regions drive the PD-like classification, providing interpretable evidence for clinical characterisation of the borderline cluster.

    \item \textbf{Test-Time Augmentation and Ensemble Inference.} Apply SpecAugment variants at inference time and aggregate predictions across augmented versions, potentially reducing the Fold~3 false positive rate by smoothing over augmentation-induced variation.

    \item \textbf{Transformer-Based Temporal Modelling.} Replace the BiLSTM with a temporal transformer block to capture longer-range dependencies and enable attention-based interpretability of the temporal patterns.

    \item \textbf{Larger Multilingual Corpora.} Extend cross-lingual evaluation to Mandarin, Arabic, and broader European language PD datasets, systematically mapping the language-dependent generalisation boundary.
\end{enumerate}

% ═══════════════════════════════════════════════════════════════
% CHAPTER 9 — CONCLUSION
% ═══════════════════════════════════════════════════════════════
\chapter{Conclusion}

This thesis has presented a novel \dualcnn{} architecture for Parkinson's Disease speech detection, combining dual-branch CNNs (ResNet50 and EfficientNetV2-S) with a custom FrequencyPriorSelfAttention module, Coordinate Attention, and a BiLSTM temporal head. The system was evaluated under speaker-independent five-fold cross-validation with certified zero-leakage guarantees, producing $N\!=\!15$ observations per dataset with Wilcoxon signed-rank statistical testing.

On the Italian PD dataset, the proposed system achieves F1\,=\,$0.964 \pm 0.048$ under speaker-independent cross-validation with certified zero leakage. No prior work on this dataset reports results under comparable evaluation rigour: the closest published figure (Aversano et al., F1\,=\,0.971) does not specify splitting strategy or leakage controls, making a direct numerical comparison uninformative. To the best of the author's knowledge, this result represents the strongest figure reported on this dataset under a fully verifiable protocol. The controlled ablation study demonstrates measurable contributions from each architectural component: Coordinate Attention ($+0.7\%$), BiLSTM ($+0.1\%$), and the EfficientNetV2-S dual branch ($+0.6\%$). Cross-lingual evaluation on the Colombian Spanish PC-GITA DDK corpus (F1\,=\,0.780) quantifies the generalisation gap across languages and establishes a baseline for future multilingual PD detection research.

The Fold~3 anomaly, a consistent performance degradation traced to approximately 21 HC subjects whose spectrograms mimic PD patterns, has implications beyond the machine learning results. This subgroup is invisible under single-split evaluation, and its existence raises questions about pre-symptomatic PD detection and whether binary clinical labels are sufficient for speech pathology datasets.

Both leakage events were caught and corrected before any final results were generated. The Strategy~A$\to$B pivot corrected a confounded experimental design before conclusions were drawn from it. These episodes share a consistent approach: methodological problems were identified, reported, and fixed rather than suppressed.

In summary, the work makes five contributions: a novel frequency-aware attention mechanism for clinical spectrograms; a certified zero-leakage evaluation pipeline; a cross-lingual PD detection evaluation under unified conditions; the identification of a borderline HC cluster not previously documented; and a documented architectural progression from baseline to the final system. The results show that under a rigorous and fully reproducible evaluation protocol, the proposed architecture produces, to the best of the author's knowledge, the strongest verifiable result on the Italian PD dataset, while being transparent about where and why the system falls short.

The core findings of this thesis have been submitted for publication as an IEEE conference paper, presenting the \dualcnn{} architecture, the zero-leakage evaluation protocol, and the cross-lingual generalisation results to the broader research community.

% ═══════════════════════════════════════════════════════════════
% CHAPTER 10 — PLANNING AND SCHEDULE
% ═══════════════════════════════════════════════════════════════
\chapter{Planning and Schedule}

\section{Project Timeline}

The project spanned two terms (October 2025 -- April 2026) with the following major phases:

\begin{table}[H]
\centering
\caption{Project Schedule: Plan vs Actual}
\label{tab:schedule}
\small
\begin{tabularx}{\textwidth}{l l X X}
\toprule
\textbf{Period} & \textbf{Phase} & \textbf{Planned} & \textbf{Actual} \\
\midrule
Oct--Dec 2025 & Term~1 & CNN baselines on ESC-50, EmoDB & Completed. CBAM instability discovered; augmentation leakage found and fixed. Scored 93/100. \\
Jan W1--2 2026 & Phase~1 & Clinical dataset integration & Completed. Italian PD, Pitt, PhysioNet integrated. Pitt bias discovered. \\
Jan W3--4 2026 & Phase~1b & Gold anchor baselines & Completed. 15-run anchor matrix. Regularisation pendulum diagnosed. \\
Feb W1--2 2026 & Phase~2 & K-fold baseline matrix & Completed. 75 runs. HybridNet dominance confirmed. Fold~3 anomaly first observed. \\
Feb W3--4 2026 & Phase~3 & Attention ablation & Extended by 1 week. Strategy~A confound required pivot to Strategy~B. \\
Mar W1 2026 & Phase~4 & PC-GITA integration & Completed. Cross-lingual DDK pipeline built. \\
Mar W1--2 2026 & Phase~5 & SOTA push & Kaggle leakage bug added 4 days. Stage~A completed; label smoothing added for Stage~B. \\
Mar W2--3 2026 & Phase~6 & Final evaluation & Completed. $N\!=\!15$ results on all datasets. Conference paper submitted. \\
Mar W3--Apr 2026 & Writing & Thesis and documentation & Completed. Thesis finalised and codebase prepared for submission. \\
\bottomrule
\end{tabularx}
\end{table}

\section{Key Deviations from Plan}

Three significant deviations from the original Term~2 plan occurred:

\textbf{Term~1 Leakage Fix (December 2025).} The discovery that augmented variants were crossing the train-validation boundary required rewriting the split logic and re-running affected experiments. This added approximately one week but established the leakage-aware methodology that became a thesis contribution.

\textbf{Phase~2 Regularisation Pendulum (February 2026).} The initial ``universal regularisation'' strategy failed catastrophically: clinical datasets overfitted while non-clinical datasets underfitted. Diagnosing and implementing the tiered protocol required three days of iterative experimentation. The tiered protocol became a documented design decision in the methodology.

\textbf{Phase~5 Kaggle Leakage Bug (March 2026).} The path-resolution leakage on Kaggle consumed four days: two for diagnosis (the bimodal F1 distribution was initially puzzling), one for the fix (hardcoded slug, RuntimeError guard), and one for re-running all affected experiments. The RuntimeError guard implemented during this fix became a permanent feature of the training pipeline and a key contribution to the evaluation methodology.

% ═══════════════════════════════════════════════════════════════
% CHAPTER 12 — PROJECT DIARY
% ═══════════════════════════════════════════════════════════════
\chapter{Project Diary}

The complete project diary is maintained in \texttt{diary.md} at the project root and version-controlled in GitLab. This chapter provides a structured summary organised by experimental phase.

\section{Term~1 (October--December 2025)}
Established the CNN baseline framework on ESC-50 and EmoDB. Key events: ResNet50 selected as primary backbone; CBAM instability discovered ($\pm 19.3\%$ on EmoDB); SE provided marginal gains; augmentation leakage identified and fixed via grouped splitting on source clip ID. Interim report submitted (93/100).

\section{Phase~1: Clinical Extension (January 2026)}
Integrated three clinical datasets: Italian PD (831 spectrograms), PhysioNet 2016 (3,153 records), and Pitt Corpus (3,836 segments with SHA-256 tracing). Discovered Pitt bias (Control Recall $<$46\%) and implemented weighted loss. Established gold anchor baselines. Encountered and resolved the regularisation pendulum: universal parameters failed, motivating the tiered protocol.

\section{Phase~2: K-Fold Baselines (February 2026)}
Transitioned to speaker-independent 5-fold CV. Executed 75 training runs. HybridNet ($\alpha$-gate fusion) dominated non-clinical tasks. ResNet50 won on Italian PD. Fold~3 anomaly first confirmed: all three models degraded to F1 $\approx$ 0.80--0.85, traced to HC subjects with PD-like spectrograms.

\section{Phase~3--4: Attention Ablation (February--March 2026)}
Designed and deployed controlled attention ablation. Identified Strategy~A confound (four variables changed simultaneously); pivoted to Strategy~B. Discovered architectural bug: 16 attention modules injected instead of 1, causing model collapse. CA identified as best single-module attention; AG rejected (8M params, no improvement).

\section{Phase~5: PC-GITA and SOTA Push (March 2026)}
Integrated PC-GITA corpus for cross-lingual evaluation. Implemented BiLSTM temporal head, SpecAugment, and dual-branch architecture. Discovered and fixed Kaggle path-resolution leakage (stale fold CSVs from prior version). Implemented RuntimeError guard. Stage~A achieved 0.960 F1 (single seed); Stage~B confirmed at 0.964 $\pm$ 0.051 ($N\!=\!15$).

\section{Phase~6: Final Evaluation and Paper (March 2026)}
Designed FrequencyPriorSelfAttention module. Completed multi-seed evaluation across all datasets. Corrected Aversano reference (year, F1, venue, authors). Removed unverified Bandini citation. Submitted conference paper. Restructured repository to RHUL FYP guidelines.

In the final submission phase, Dr.\ Li's feedback was fully implemented: PC-GITA counts and result tables were corrected, stronger baseline comparisons were added, the final publication figures were generated and inserted, and the repository was restructured to conform to RHUL FYP submission guidelines. Phase~6 was also hardened at the implementation level through fixes to notebook generation, variable-scoping issues, dataset packaging problems, and additional runtime checks to ensure reproducible execution.

% ═══════════════════════════════════════════════════════════════
% APPENDICES
% ═══════════════════════════════════════════════════════════════
\appendix

\chapter{README and Directory Structure}

The project README (\texttt{README.md}) provides complete setup, execution, and reproduction instructions. The repository structure is:

\begin{verbatim}
product/
├── audio_preprocessing/     Spectrogram generation pipeline
│   ├── data/                Raw audio datasets (excluded)
│   ├── src/                 Generation scripts per dataset
│   └── outputs/             Generated spectrogram PNGs
├── models/                  All model architectures
│   ├── dual_cnn_sa_lstm.py  Novel architecture
│   ├── freq_prior_attention.py  FP-SA module
│   ├── model_builder.py     Factory pattern
│   └── [other modules]
├── training/
│   └── train_unified.py     Dataset-agnostic training engine
├── artifacts/
│   ├── splits/              All fold CSVs (50 files)
│   └── runs/                Results, checkpoints, reports
├── notebooks/               Kaggle execution notebooks
└── scripts/                 Utilities and verification tools

documents/                   Reports and paper
.gitlab-ci.yml               CI/CD pipeline
diary.md                     Complete project diary
\end{verbatim}

\chapter{Installation and Execution Instructions}

\section{Environment Setup}
\begin{verbatim}
python -m venv .venv
.venv\Scripts\activate          # Windows
pip install -U pip
pip install numpy scipy librosa soundfile matplotlib \
    pandas scikit-learn torch torchvision tensorboard pillow
\end{verbatim}

\section{Dataset Acquisition}
Due to privacy and ethics agreements, the raw audio datasets are not included in the repository and must be acquired independently.

\begin{itemize}
    \item \textbf{Italian PD:} Available publicly on IEEE Dataport (\url{https://ieee-dataport.org/open-access/italian-parkinsons-voice-and-speech}). Place extracted files in \texttt{product/audio\_preprocessing/data/Italian Parkinson's Voice and speech/}
    \item \textbf{PC-GITA:} Requires special access permission. Request directly from Prof. Juan Rafael Orozco-Arroyave. Once granted, place extracted files in \texttt{product/audio\_preprocessing/data/PC-GITA/}
    \item \textbf{ESC-50:} Available publicly on GitHub at \url{https://github.com/karolpiczak/ESC-50}. Place extracted files in \texttt{product/audio\_preprocessing/data/ESC-50/}
    \item \textbf{EmoDB:} Available publicly from TU Berlin at \url{http://emodb.bilderbar.info/download/}. Place extracted files in \texttt{product/audio\_preprocessing/data/EmoDB-wav/}
\end{itemize}

\section{Spectrogram Generation}
\begin{verbatim}
python product/audio_preprocessing/src/generate_spectrograms.py
python product/audio_preprocessing/src/generate_spectrograms_emodb.py
python product/audio_preprocessing/src/generate_spectrograms_italian.py
python product/audio_preprocessing/src/generate_spectrograms_pcgita.py
\end{verbatim}

\section{K-Fold Split Generation}
\begin{verbatim}
python scripts/generate_kfold_splits.py --n_folds 5 --seed 42
\end{verbatim}

\section{Training}
\begin{verbatim}
# Single fold
python product/training/train_unified.py \
    --dataset italian_pd --model_type dual_cnn_sa_lstm \
    --fold 0 --seed 42 --epochs 30 --spec_augment \
    --label_smoothing 0.05

# Full 75-run matrix
scripts\run_kfold_experiments.bat
\end{verbatim}

\section{Result Aggregation}
\begin{verbatim}
python scripts/aggregate_kfold_results.py --n_folds 5
\end{verbatim}

\section{Expected System Output and Screenshots}
\label{sec:screenshots}

This section documents the expected terminal output and generated artefacts when the system executes correctly. The marker should compare their output against these reference examples to verify successful execution.

\subsection*{Leakage Verification Output}

Running \texttt{verify\_leakage.py} should produce output in the following form:
\begin{verbatim}
Checking 50 fold CSV files (5 datasets x 5 folds x 2 splits)...
[italian_pd] Fold 0: PASS | Fold 1: PASS | Fold 2: PASS |
             Fold 3: PASS | Fold 4: PASS
[pc_gita]    Fold 0: PASS | ... | Fold 4: PASS
...
------------------------------------------------------------
Result: 0 violations across 50 fold files.
All subject-level leakage checks PASSED.
\end{verbatim}

\subsection*{Training Run Output (Single Fold)}

A correctly executing training run produces console output in the following format:
\begin{verbatim}
[Leakage Guard] Checking fold 0 for subject overlap...
[Leakage Guard] PASSED -- 0 overlapping subjects detected.
Epoch  1/30 | Train Loss: 0.6821 | Val F1: 0.721
Epoch 10/30 | Train Loss: 0.2876 | Val F1: 0.941
Epoch 30/30 | Train Loss: 0.1543 | Val F1: 0.984
[Checkpoint] Best model saved at epoch 28 (F1: 0.991)
[Artefacts]  Saved: summary.json, confusion_matrix.png,
             best_classification_report.json, best_model.pt
\end{verbatim}

\subsection*{Aggregated Results Output}

After completing all 15 runs on Italian PD, \texttt{aggregate\_kfold\_results.py} produces a summary. The Italian PD aggregate should read $0.964 \pm 0.048$ for \texttt{dual\_cnn\_sa\_lstm}.

\subsection*{Verification and Output Artefacts}

Each training run produces three artefacts saved to \texttt{product/artifacts/runs/<dataset>/<run\_name>/}: a \texttt{confusion\_matrix.png} visualising per-class prediction accuracy, a \texttt{classification\_report.json} containing precision, recall, F1, and support per class, and a \texttt{model\_best.pt} checkpoint storing the weights at peak validation F1. The leakage guard executes before any training begins and prints a \texttt{PASS} confirmation to stdout; if any subject overlap is detected, a \texttt{RuntimeError} is raised and execution halts immediately, preventing any result from being recorded under compromised conditions.

A successful run on Italian PD (\texttt{dual\_cnn\_sa\_lstm}, fold~0, seed~42) converges within 30 epochs, with validation F1 typically rising from approximately 0.71 at epoch~1 to 0.95--0.96 by epoch~30. The confusion matrix for fold~0 exhibits near-perfect diagonal structure, reflecting clean speaker-independent separation. Fold~3, by contrast, shows a characteristic pattern of elevated healthy control false positives --- the visual signature of the borderline HC cluster described in Section~\ref{subsec:fold3} --- which persists consistently across all architectures and all random seeds evaluated in this project.

\chapter{Additional Results Tables}

\section{Phase~5 Stage~A: Italian PD Full Fold Breakdown (seed 42)}

\begin{table}[H]
\centering
\caption{Stage~A Italian PD Per-Fold F1 (seed 42)}
\begin{tabular}{l c c c c c c}
\toprule
\textbf{Model} & \textbf{F0} & \textbf{F1} & \textbf{F2} & \textbf{F3} & \textbf{F4} & \textbf{Mean $\pm$ Std} \\
\midrule
ResNet50 & 0.984 & 0.987 & 0.991 & 0.865 & 0.923 & 0.950 $\pm$ 0.049 \\
+CA & 0.984 & 0.987 & 0.991 & 0.885 & 0.936 & 0.957 $\pm$ 0.041 \\
+CA+BiLSTM & 1.000 & 1.000 & 1.000 & 0.858 & 0.932 & 0.958 $\pm$ 0.058 \\
\texttt{dual\_cnn\_lstm} & 1.000 & 1.000 & 1.000 & 0.878 & 0.941 & 0.964 $\pm$ 0.051 \\
\bottomrule
\end{tabular}
\end{table}

\section{Phase~4: PC-GITA Cross-Lingual Fold Breakdown}

\begin{table}[H]
\centering
\caption{PC-GITA Per-Fold Baseline Evaluation (\dualcnn{})}
\begin{tabular}{l c c c c c c}
\toprule
\textbf{Model} & \textbf{F0} & \textbf{F1} & \textbf{F2} & \textbf{F3} & \textbf{F4} & \textbf{Mean $\pm$ Std} \\
\midrule
\dualcnn{} & 0.754 & 0.800 & 0.830 & 0.691 & 0.824 & 0.780 $\pm$ 0.056 \\
\bottomrule
\end{tabular}
\end{table}

\section{Phase~2: Italian PD Fold~3 Anomaly Detail}

\begin{table}[H]
\centering
\caption{Phase~2 Italian PD Fold~3 Performance Pattern}
\begin{tabular}{l c c}
\toprule
\textbf{Metric} & \textbf{Folds 0--2, 4 (avg)} & \textbf{Fold~3} \\
\midrule
PD Recall & $\sim$0.98 & 1.000 \\
HC Recall & $\sim$0.90 & $\sim$0.75 \\
Macro F1 & $\sim$0.95 & $\sim$0.85 \\
\bottomrule
\end{tabular}
\end{table}

\section{Fold~3 Anomaly: Seed-Invariance Confirmation}

Table~\ref{tab:fold3_seed_consistency} confirms that the Fold~3 performance degradation is consistent across all three random seeds in the final \dualcnn{} evaluation. If the anomaly were an artefact of a particular weight initialisation, it would not persist across seeds. In all three seeds, Fold~3 produces the lowest F1 score of any fold, confirming the anomaly as a structural property of the Fold~3 validation set composition rather than an initialisation artefact.

\begin{table}[H]
\centering
\caption{Italian PD Per-Fold F1 Across All Seeds (\dualcnn{}). Fold~3 is consistently the lowest across all seeds.}
\label{tab:fold3_seed_consistency}
\begin{tabular}{c c c c c c c}
\toprule
\textbf{Seed} & \textbf{Fold 0} & \textbf{Fold 1} & \textbf{Fold 2} & \textbf{Fold 3} & \textbf{Fold 4} & \textbf{Mean} \\
\midrule
42  & 0.994 & 1.000 & 0.991 & 0.865 & 0.954 & 0.961 \\
123 & 1.000 & 1.000 & 1.000 & 0.898 & 0.945 & 0.969 \\
999 & 0.994 & 1.000 & 1.000 & 0.884 & 0.937 & 0.963 \\
\midrule
\textbf{Fold mean} & 0.996 & 1.000 & 0.997 & \textbf{0.882} & 0.945 & \textbf{0.964} \\
\bottomrule
\end{tabular}
\end{table}

\chapter{User Manual}
\label{app:usermanual}

This appendix describes how to use the \dualcnn{} system for PD speech detection. It is written for a technical reader who wants to reproduce the reported results or apply the trained pipeline to new audio data. For environment setup and dataset acquisition, see Appendix~B.

\section{Overview}

The system takes raw speech audio files as input and produces binary classification outputs (PD or Healthy Control) alongside confidence scores. The pipeline has two stages: (1)~spectrogram generation, which converts audio to $224 \times 224$ log-Mel spectrogram PNG images; and (2)~model inference, which classifies spectrograms using the trained \dualcnn{} model.

\section{Reproducing the Reported Results}

The reported results (Italian PD: $\text{F1} = 0.964 \pm 0.048$, $N\!=\!15$) are reproducible by executing the full K-fold evaluation pipeline. All fold split CSVs and trained model checkpoints are included in the submission under \texttt{product/artifacts/}.

\textbf{Step 1: Verify leakage integrity.}
\begin{verbatim}
python scripts/verify_leakage.py
\end{verbatim}
Expected output: \texttt{0 violations across 50 fold files. All checks PASSED.}

\textbf{Step 2: Run a single fold end-to-end.}
\begin{verbatim}
python product/training/train_unified.py \
    --dataset italian_pd \
    --model_type dual_cnn_sa_lstm \
    --fold 0 --seed 42 --epochs 30 \
    --spec_augment --label_smoothing 0.05
\end{verbatim}

\textbf{Step 3: Run the full 15-observation evaluation.}
\begin{verbatim}
scripts\run_kfold_experiments.bat
python scripts/aggregate_kfold_results.py --n_folds 5
\end{verbatim}
Expected: \texttt{aggregated\_results.json} with Italian PD overall mean F1\,=\,$0.964 \pm 0.048$.

\section{Output Files Reference}

Each completed training run saves the following artefacts to \texttt{product/artifacts/runs/\{dataset\}/\{model\}/seed\{N\}/fold\{K\}/}:

\begin{itemize}
    \item \textbf{\texttt{summary.json}} --- Contains \texttt{best\_val\_f1}, \texttt{best\_val\_accuracy}, \texttt{best\_val\_auc}, \texttt{best\_epoch}, and the full hyperparameter configuration. The \texttt{best\_val\_f1} field is the primary metric used throughout this thesis.
    \item \textbf{\texttt{best\_classification\_report.json}} --- Scikit-learn classification report in JSON format. The \texttt{recall} of the \texttt{HC} entry is Control Recall, the key secondary metric for clinical safety. On Fold~3, this value will be lower than on other folds, reflecting the borderline HC cluster (Section~\ref{subsec:fold3}).
    \item \textbf{\texttt{confusion\_matrix.png}} --- $2 \times 2$ confusion matrix. Rows are true class labels; columns are predicted labels. For Fold~3, the HC row will show elevated off-diagonal counts.
    \item \textbf{\texttt{best\_model.pt}} --- PyTorch state dictionary for the checkpoint achieving the best validation F1. Load with \texttt{torch.load()} and the corresponding \texttt{model\_builder} factory call.
\end{itemize}

\section{Applying the Model to New Audio}

To classify new recordings using a trained checkpoint:
\begin{verbatim}
python scripts/infer.py \
    --audio_path /path/to/recording.wav \
    --model_checkpoint product/artifacts/runs/italian_pd/ \
        dual_cnn_sa_lstm/seed42/fold0/best_model.pt \
    --model_type dual_cnn_sa_lstm
\end{verbatim}
Output: \texttt{\{"prediction": "PD"|"HC", "confidence": 0.0--1.0\}}.

\textbf{Important:} This system is a research prototype for academic evaluation only. It has not undergone clinical validation and must not be used for diagnostic purposes.

\section{Verifying Data Integrity}

SHA-256 hashes embedded in spectrogram filenames (e.g., \texttt{157\_2c87d4e1a0b5\_015\_orig.png}) can be verified against source recordings:
\begin{verbatim}
python scripts/verify_hashes.py --dataset italian_pd
\end{verbatim}
Expected output: \texttt{All 831 spectrograms verified. Zero hash mismatches.}

\chapter{AI Use Statement}

\section*{Tools Used}

\begin{itemize}
    \item \textbf{Claude Code} --- Anthropic, \url{https://claude.ai}.
    \item \textbf{Antigravity IDE} --- Google DeepMind, \url{https://deepmind.google}.
\end{itemize}

\section*{Description of Use}

The above tools were used solely for supportive purposes, including:
\begin{itemize}
    \item Clarifying technical concepts to aid my own understanding (e.g., advanced convolutional neural network architectures, transfer learning, and attention mechanisms such as Squeeze-and-Excitation, CBAM, and Coordinate Attention).
    \item Providing code execution assistance, AI-assisted code completion, and debugging support during the development of the training pipeline.
    \item Improving clarity, structure, and grammar of text written by me, without altering technical meaning or content.
    \item Assisting with high-level planning and organisation of the report (e.g., structuring sections and timelines).
    \item Generating illustrative system diagrams and results plots (including the dual-branch architecture diagram, ablation study charts, and fold-level anomaly visualisations) based strictly on specifications, data, and workflows defined by me. These figures were used solely for visual communication and presentation clarity. All underlying architectural decisions, system design choices, and implementations are entirely my own.
\end{itemize}

All model architectures (including the novel \dualcnn{} and FrequencyPriorSelfAttention module), training pipelines, code implementations, experimental design, hyperparameter choices, evaluation procedures, results, and analysis were developed and written by me.

Where external resources or publicly available implementations were consulted, they are explicitly acknowledged and referenced in the report in accordance with standard academic practice.

No AI-generated text or content has been included verbatim without review, modification, and integration into my own work.

I confirm that the use of generative AI tools in this project complies fully with the department's policy on the responsible use of AI, and that this submission represents my own independent work.

\vspace{1.5cm}

\noindent Signed: \includegraphics[height=1.2cm]{figures/signature.png} \\
\noindent \textbf{Devansh Dev} \\
BSc Computer Science with Artificial Intelligence \\
Royal Holloway, University of London

%%%% BIBLIOGRAPHY
\newpage
\begin{thebibliography}{99}
\addcontentsline{toc}{chapter}{Bibliography}

\bibitem{dorsey2018} E.\ R.\ Dorsey, T.\ Sherer, M.\ S.\ Okun, and B.\ R.\ Bloem, ``The emerging evidence of the Parkinson pandemic,'' \emph{J.\ Parkinson's Dis.}, vol.~8, no.~S1, pp.~S3--S8, 2018.

\bibitem{ho1999} A.\ K.\ Ho, R.\ Iansek, C.\ Marigliani, J.\ L.\ Bradshaw, and S.\ Gates, ``Speech impairment in a large sample of patients with Parkinson's disease,'' \emph{Behav.\ Neurol.}, vol.~11, no.~3, pp.~131--137, 1999.

\bibitem{duffy2013} J.\ R.\ Duffy, \emph{Motor Speech Disorders: Substrates, Differential Diagnosis, and Management}, 3rd~ed.\ St.\ Louis, MO: Elsevier Mosby, 2013.

\bibitem{sakar2013} B.\ E.\ Sakar, M.\ E.\ Isenkul, C.\ O.\ Sakar, A.\ Sertbas, F.\ Gurgen, S.\ Delil, H.\ Apaydin, and O.\ Kursun, ``Collection and analysis of a Parkinson speech dataset with multiple types of sound recordings,'' \emph{IEEE J.\ Biomed.\ Health Inform.}, vol.~17, no.~4, pp.~828--834, 2013.

\bibitem{tsanas2012} A.\ Tsanas, M.\ A.\ Little, C.\ Fox, and L.\ O.\ Ramig, ``Objective automatic assessment of rehabilitative speech treatment in Parkinson's disease,'' \emph{IEEE Trans.\ Neural Syst.\ Rehabil.\ Eng.}, vol.~22, no.~1, pp.~181--190, 2014.

\bibitem{he2016} K.\ He, X.\ Zhang, S.\ Ren, and J.\ Sun, ``Deep residual learning for image recognition,'' in \emph{Proc.\ IEEE CVPR}, 2016, pp.~770--778.

\bibitem{tan2021} M.\ Tan and Q.\ V.\ Le, ``EfficientNetV2: Smaller models and faster training,'' in \emph{Proc.\ ICML}, 2021, pp.~10096--10106.

\bibitem{aversano2024} L.\ Aversano, M.\ L.\ Bernardi, M.\ Cimitile, R.\ Pecori, A.\ Picariello, G.\ Sperl\`{i}, and A.\ Testa, ``A machine learning approach for early detection of Parkinson's disease using acoustic traces,'' in \emph{Proc.\ IEEE EAIS}, 2022, pp.~1--8.

\bibitem{orozco2014} J.\ R.\ Orozco-Arroyave, J.\ D.\ Arias-Londo\~{n}o, J.\ F.\ Vargas-Bonilla, M.\ C.\ Gonzalez-Ratvia, and E.\ N\"{o}th, ``New Spanish speech corpus database for the analysis of people suffering from Parkinson's disease,'' in \emph{Proc.\ LREC}, 2014, pp.~342--347.

\bibitem{hu2018} J.\ Hu, L.\ Shen, and G.\ Sun, ``Squeeze-and-excitation networks,'' in \emph{Proc.\ IEEE CVPR}, 2018, pp.~7132--7141.

\bibitem{woo2018} S.\ Woo, J.\ Park, J.-Y.\ Lee, and I.\ S.\ Kweon, ``CBAM: Convolutional block attention module,'' in \emph{Proc.\ ECCV}, 2018, pp.~3--19.

\bibitem{hou2021} Q.\ Hou, D.\ Zhou, and J.\ Feng, ``Coordinate attention for efficient mobile network design,'' in \emph{Proc.\ IEEE/CVF CVPR}, 2021, pp.~13713--13722.

\bibitem{rusz2011} J.\ Rusz, R.\ Cmejla, H.\ Ruzickova, and E.\ Ruzicka, ``Quantitative acoustic measurements for characterization of speech and voice disorders in early untreated Parkinson's disease,'' \emph{J.\ Acoust.\ Soc.\ Am.}, vol.~129, no.~1, pp.~350--367, 2011.

\bibitem{graves2005} A.\ Graves and J.\ Schmidhuber, ``Framewise phoneme classification with bidirectional LSTM and other neural network architectures,'' \emph{Neural Netw.}, vol.~18, no.~5--6, pp.~602--610, 2005.

\bibitem{hochreiter1997} S.\ Hochreiter and J.\ Schmidhuber, ``Long short-term memory,'' \emph{Neural Comput.}, vol.~9, no.~8, pp.~1735--1780, 1997.

\bibitem{dimauro2019} G.\ Dimauro, V.\ Di Nicola, V.\ Bevilacqua, D.\ Caivano, and F.\ Girardi, ``Italian Parkinson's voice and speech,'' \emph{IEEE Dataport}, 2019.

\bibitem{orozco2016} J.\ R.\ Orozco-Arroyave, F.\ Hoenig, J.\ D.\ Arias-Londo\~{n}o, J.\ F.\ Vargas-Bonilla, K.\ Daqrouq, S.\ Hadjitodorov, and E.\ N\"{o}th, ``Automatic detection of Parkinson's disease in running speech spoken in three different languages,'' \emph{J.\ Acoust.\ Soc.\ Am.}, vol.~139, no.~1, pp.~481--500, 2016.

\bibitem{vaswani2017} A.\ Vaswani, N.\ Shazeer, N.\ Parmar, J.\ Uszkoreit, L.\ Jones, A.\ N.\ Gomez, L.\ Kaiser, and I.\ Polosukhin, ``Attention is all you need,'' in \emph{Proc.\ NeurIPS}, 2017.

\bibitem{piczak2015} K.\ J.\ Piczak, ``ESC: Dataset for environmental sound classification,'' in \emph{Proc.\ ACM Multimedia}, 2015.

\bibitem{burkhardt2005} F.\ Burkhardt, A.\ Paeschke, M.\ Rolfes, W.\ Sendlmeier, and B.\ Weiss, ``A database of German emotional speech,'' in \emph{Proc.\ Interspeech}, 2005.

\bibitem{park2019} D.\ S.\ Park, W.\ Chan, Y.\ Zhang, C.-C.\ Chiu, B.\ Zoph, E.\ D.\ Cubuk, and Q.\ V.\ Le, ``SpecAugment: A simple data augmentation method for automatic speech recognition,'' in \emph{Proc.\ Interspeech}, 2019.

\bibitem{sklearn2011} F.\ Pedregosa et al., ``Scikit-learn: Machine learning in Python,'' \emph{J.\ Mach.\ Learn.\ Res.}, vol.~12, pp.~2825--2830, 2011.

\end{thebibliography}
\label{endpage}

\end{document}
